\chapter{DX-031 Standard Bijou Blocks}\label{record-128}

\noindent{\footnotesize\color{BookMuted}\path{docs/design/DX-031-standard-bijou-blocks.md}\par}\smallskip

\begin{RecordMeta}
title: DX-031 Standard Bijou Blocks
legend: DX
lane: design

\end{RecordMeta}

\emph{Cycle for declaring reusable app-building blocks that compose Bijou components into standard, inspectable surface patterns}

Legend:

\begin{itemize}
\tightlist
\item
  \href{https://github.com/flyingrobots/bijou/blob/1853ad902f6bf219326409e7f743a5a1d6517dbc/docs/legends/DX-developer-experience.md}{DX - Developer Experience}
\end{itemize}

Depends on:

\begin{itemize}
\tightlist
\item
  \href{https://github.com/flyingrobots/bijou/blob/cbd70612140e7a450d764b3c7be1c6af878b99bc/docs/design/DX-025-component-metadata-contract.md}{DX-025 - Component Metadata Contract}
\item
  \href{https://github.com/flyingrobots/bijou/blob/cbd70612140e7a450d764b3c7be1c6af878b99bc/docs/design/DX-031F-block-tree-rendering.md}{DX-031F - Block Tree Rendering}
\item
  \href{https://github.com/flyingrobots/bijou/blob/cbd70612140e7a450d764b3c7be1c6af878b99bc/docs/design/DF-028-story-capture-matrix.md}{DF-028 - Story Capture Matrix}
\item
  \href{https://github.com/flyingrobots/bijou/blob/cbd70612140e7a450d764b3c7be1c6af878b99bc/docs/design/DF-029-fixture-to-docs-promotion-path.md}{DF-029 - Fixture-to-Docs Promotion Path}
\item
  \href{https://github.com/flyingrobots/bijou/blob/cbd70612140e7a450d764b3c7be1c6af878b99bc/docs/design/RE-035-mandatory-layout-envelope-and-constraint-negotiation.md}{RE-035}
\end{itemize}

\section{Why This Cycle Exists}\label{record-128-why-this-cycle-exists}

Bijou has many capable components, but app builders still need to assemble common surfaces by hand: shells, readers, inspectors, command palettes, settings, modal stacks, notifications, and authentication flows.

Blocks standardize those common compositions without turning them into new primitives.

A block is a named composition of Bijou components with a documented layout, focus, input, mode-lowering, and semantic contract. Blocks are the layer a builder reaches for when they want to build an app quickly without re-deciding the same shell and workflow patterns every time.

\section{Boundary}\label{record-128-boundary}

\begin{itemize}
\tightlist
\item
  Blocks compose primitives.
\item
  Blocks expose slots and semantic facts.
\item
  Blocks lower across interactive, static, pipe, and accessible modes.
\item
  Blocks do not mutate layout directly.
\item
  Blocks do not replace RE layout primitives.
\item
  Blocks do not become a hidden window manager.
\end{itemize}

\section{Relationship To RE-035}\label{record-128-relationship-to-re-035}

Blocks depend on layout truth, but they do not define it.

Each block should compile into ordinary visible nodes with mandatory layout envelopes. The block owns composition, slot names, default behavior, and semantic facts. The layout engine owns constraints, preferences, assigned rectangles, and explanations.

That keeps the boundary simple:

\begin{itemize}
\tightlist
\item
  blocks choose what app surface exists
\item
  layout chooses where each visible node fits
\item
  focus and input consume retained layout facts
\item
  renderers paint assigned rectangles
\item
  modes lower semantic block facts
\end{itemize}

\section{Runtime Truth Vs App Truth}\label{record-128-runtime-truth-vs-app-truth}

Blocks must preserve a hard boundary between primitive runtime truth and semantic app truth.

Primitive runtime truth belongs to RE and adjacent runtime layers:

\begin{itemize}
\tightlist
\item
  layout envelope
\item
  constraints
\item
  preferences
\item
  assigned rectangles
\item
  fit policy
\item
  focus regions
\item
  hit regions
\item
  viewport facts
\item
  layer order
\item
  command and effect facts
\end{itemize}

Semantic app truth belongs to blocks and app code:

\begin{itemize}
\tightlist
\item
  \texttt{R\allowbreak{}e\allowbreak{}a\allowbreak{}d\allowbreak{}e\allowbreak{}r\allowbreak{}S\allowbreak{}u\allowbreak{}r\allowbreak{}f\allowbreak{}a\allowbreak{}c\allowbreak{}e}
\item
  \texttt{P\allowbreak{}r\allowbreak{}o\allowbreak{}d\allowbreak{}u\allowbreak{}c\allowbreak{}t\allowbreak{}D\allowbreak{}e\allowbreak{}t\allowbreak{}a\allowbreak{}i\allowbreak{}l\allowbreak{}s}
\item
  \texttt{C\allowbreak{}h\allowbreak{}e\allowbreak{}c\allowbreak{}k\allowbreak{}o\allowbreak{}u\allowbreak{}t\allowbreak{}F\allowbreak{}l\allowbreak{}o\allowbreak{}w}
\item
  \texttt{D\allowbreak{}a\allowbreak{}s\allowbreak{}h\allowbreak{}b\allowbreak{}o\allowbreak{}a\allowbreak{}r\allowbreak{}d\allowbreak{}S\allowbreak{}h\allowbreak{}e\allowbreak{}l\allowbreak{}l}
\item
  \texttt{C\allowbreak{}o\allowbreak{}m\allowbreak{}m\allowbreak{}e\allowbreak{}n\allowbreak{}t\allowbreak{}s\allowbreak{}T\allowbreak{}h\allowbreak{}r\allowbreak{}e\allowbreak{}a\allowbreak{}d}
\item
  \texttt{T\allowbreak{}e\allowbreak{}r\allowbreak{}m\allowbreak{}i\allowbreak{}n\allowbreak{}a\allowbreak{}l\allowbreak{}P\allowbreak{}a\allowbreak{}n\allowbreak{}e\allowbreak{}l}
\item
  domain entities, statuses, actions, and relationships
\end{itemize}

The dependency direction only goes one way:

\begin{itemize}
\tightlist
\item
  semantic blocks compile into primitive runtime truth
\item
  runtime primitives do not learn product, commerce, docs, or chat semantics
\item
  app semantics do not mutate runtime geometry directly
\item
  inspectors may relate both layers, but must label which layer they explain
\end{itemize}

This prevents business semantics from leaking into the layout engine and keeps runtime primitives from becoming app-shaped.

\section{Human Users / Jobs / Hill}\label{record-128-human-users-jobs-hill}

\subsection{Primary Human Users}\label{record-128-primary-human-users}

\begin{itemize}
\tightlist
\item
  App builders assembling Bijou tools from common surface patterns.
\item
  DOGFOOD maintainers turning docs, inspectors, settings, and command surfaces into reusable examples.
\item
  Component authors who need a standard target for real-world composition stories.
\end{itemize}

\subsection{Human Jobs}\label{record-128-human-jobs}

\begin{enumerate}
\def\labelenumi{\arabic{enumi}.}
\tightlist
\item
  Pick a block for a common app surface instead of starting from blank layout.
\item
  Fill named slots with app-owned content.
\item
  Trust the block's default focus, input, responsive, and mode-lowering behavior.
\item
  Inspect how a block composes lower-level components.
\end{enumerate}

\subsection{Human Hill}\label{record-128-human-hill}

A builder can create a useful Bijou app shell with navigation, content, settings, help, notifications, and modal behavior by composing documented blocks instead of hand-rolling every surface.

\section{Agent Users / Jobs / Hill}\label{record-128-agent-users-jobs-hill}

\subsection{Primary Agent Users}\label{record-128-primary-agent-users}

\begin{itemize}
\tightlist
\item
  Agents generating Bijou app scaffolds.
\item
  Agents upgrading examples into DOGFOOD stories.
\item
  Agents explaining how an app surface is composed.
\end{itemize}

\subsection{Agent Jobs}\label{record-128-agent-jobs}

\begin{enumerate}
\def\labelenumi{\arabic{enumi}.}
\tightlist
\item
  Select an appropriate block from a stable catalog.
\item
  Fill required slots and avoid unsupported private layout behavior.
\item
  Generate tests against block contracts instead of snapshot-only output.
\item
  Preserve semantic facts when lowering to pipe and accessible modes.
\end{enumerate}

\subsection{Agent Hill}\label{record-128-agent-hill}

An agent can inspect a block definition, identify its slots and mode behavior, and generate type-safe app code that uses the block without inventing new layout or input conventions.

\section{Accessibility / Assistive Posture}\label{record-128-accessibility-assistive-posture}

Blocks must preserve semantic meaning when visual chrome lowers away.

Every standard block should identify:

\begin{itemize}
\tightlist
\item
  named regions
\item
  active item or selected entity, when relevant
\item
  available actions
\item
  validation or warning facts
\item
  modal or focus-trap state, when relevant
\item
  whether hidden visual content remains available in accessible mode
\end{itemize}

Auth blocks must never expose secrets in accessible, pipe, static, or diagnostic output.

\section{Localization / Directionality Posture}\label{record-128-localization-directionality-posture}

Blocks should describe layout in logical terms and inherit direction from the host context unless a block explicitly owns a local direction boundary.

The first catalog should avoid physical left/right naming in slot contracts. For example:

\begin{itemize}
\tightlist
\item
  use \texttt{n\allowbreak{}a\allowbreak{}v\allowbreak{}i\allowbreak{}g\allowbreak{}a\allowbreak{}t\allowbreak{}i\allowbreak{}o\allowbreak{}n} instead of \texttt{l\allowbreak{}e\allowbreak{}f\allowbreak{}t\allowbreak{}N\allowbreak{}a\allowbreak{}v}
\item
  use \texttt{i\allowbreak{}n\allowbreak{}s\allowbreak{}p\allowbreak{}e\allowbreak{}c\allowbreak{}t\allowbreak{}o\allowbreak{}r} instead of \texttt{r\allowbreak{}i\allowbreak{}g\allowbreak{}h\allowbreak{}t\allowbreak{}P\allowbreak{}a\allowbreak{}n\allowbreak{}e}
\item
  use \texttt{p\allowbreak{}r\allowbreak{}i\allowbreak{}m\allowbreak{}a\allowbreak{}r\allowbreak{}y\allowbreak{}A\allowbreak{}c\allowbreak{}t\allowbreak{}i\allowbreak{}o\allowbreak{}n} and \texttt{s\allowbreak{}e\allowbreak{}c\allowbreak{}o\allowbreak{}n\allowbreak{}d\allowbreak{}a\allowbreak{}r\allowbreak{}y\allowbreak{}A\allowbreak{}c\allowbreak{}t\allowbreak{}i\allowbreak{}o\allowbreak{}n\allowbreak{}s} instead of button order names
\end{itemize}

Text, bidi, and locale-specific formatting stay outside this first blocks cycle, but block slots must not make later localization harder.

\section{Agent Inspectability / Explainability Posture}\label{record-128-agent-inspectability-explainability-posture}

Blocks should be inspectable as data before rendering. An agent should be able to answer:

\begin{itemize}
\tightlist
\item
  which block is this?
\item
  which slots are required?
\item
  which slots are filled?
\item
  which lower-level components does it compose?
\item
  what semantic facts does it expose?
\item
  how does it lower by mode?
\item
  what are the expected focus and action surfaces?
\end{itemize}

The first implementation should prefer metadata and story captures over snapshot-only proof.

\section{Linked Invariants}\label{record-128-linked-invariants}

\begin{itemize}
\tightlist
\item
  \href{https://github.com/flyingrobots/bijou/blob/cbd70612140e7a450d764b3c7be1c6af878b99bc/docs/invariants/docs-are-the-demo.md}{Docs Are the Demo}
\item
  \href{https://github.com/flyingrobots/bijou/blob/cbd70612140e7a450d764b3c7be1c6af878b99bc/docs/invariants/runtime-truth-wins.md}{Runtime Truth Wins}
\item
  \href{https://github.com/flyingrobots/bijou/blob/cbd70612140e7a450d764b3c7be1c6af878b99bc/docs/invariants/layout-owns-interaction-geometry.md}{Layout Owns Interaction Geometry}
\item
  \href{https://github.com/flyingrobots/bijou/blob/cbd70612140e7a450d764b3c7be1c6af878b99bc/docs/invariants/buffer-holds-facts.md}{The Buffer Holds Facts}
\item
  \href{https://github.com/flyingrobots/bijou/blob/cbd70612140e7a450d764b3c7be1c6af878b99bc/docs/invariants/commands-change-state-effects-do-not.md}{Commands Change State, Effects Do Not}
\item
  \href{https://github.com/flyingrobots/bijou/blob/cbd70612140e7a450d764b3c7be1c6af878b99bc/docs/invariants/tests-are-the-spec.md}{Tests Are the Spec}
\end{itemize}

\section{Block Contract}\label{record-128-block-contract}

Every standard block should declare:

\begin{itemize}
\tightlist
\item
  purpose
\item
  required slots
\item
  optional slots
\item
  variants
\item
  configurable options
\item
  composed components
\item
  layout posture
\item
  focus posture
\item
  input posture
\item
  responsive posture
\item
  mode-lowering posture
\item
  accessibility facts
\item
  examples and story ids
\item
  tests that define the contract
\end{itemize}

The eventual metadata can align with \texttt{d\allowbreak{}e\allowbreak{}f\allowbreak{}i\allowbreak{}n\allowbreak{}e\allowbreak{}C\allowbreak{}o\allowbreak{}m\allowbreak{}p\allowbreak{}o\allowbreak{}n\allowbreak{}e\allowbreak{}n\allowbreak{}t\allowbreak{}M\allowbreak{}e\allowbreak{}t\allowbreak{}a\allowbreak{}d\allowbreak{}a\allowbreak{}t\allowbreak{}a\allowbreak{}(\allowbreak{})} and story capture records, but this cycle starts by standardizing the human-readable catalog.

The entity relationship is:

\begin{center}
\BookDiagram[.9\linewidth]{%
\begin{tikzpicture}[classbox/.style={draw=BookHeading,fill=BookPanel,line width=.6pt,rounded corners=1.5pt,text width=37mm,align=left,inner sep=5pt,font=\scriptsize},rel/.style={-{Latex[length=2mm]},draw=BookHeading,line width=.6pt}]
\node[classbox] (BlockPackageManifest) at (0,0) {\textbf{BlockPackageManifest}\\[-2pt]\rule{\linewidth}{.3pt}\\packageName\\version\\bijouPeerRange};
\node[classbox] (BlockDefinition) at (4.7,0) {\textbf{BlockDefinition}\\[-2pt]\rule{\linewidth}{.3pt}\\metadata\\render};
\node[classbox] (BlockMetadata) at (9.4,0) {\textbf{BlockMetadata}\\[-2pt]\rule{\linewidth}{.3pt}\\id\\name\\purpose};
\node[classbox] (SchemaBoundBlock) at (4.7,-3.5) {\textbf{SchemaBoundBlock}\\[-2pt]\rule{\linewidth}{.3pt}\\schema\\bind};
\node[classbox] (BlockSchemaAdapter) at (0,-3.5) {\textbf{BlockSchemaAdapter}\\[-2pt]\rule{\linewidth}{.3pt}\\id\\parse};
\draw[rel,{Diamond[length=2mm]}-] (BlockPackageManifest) to node[midway,fill=white,inner sep=2pt,font=\scriptsize,text=BookMuted,align=center] {exports / 1 : 1..*} (BlockDefinition);
\draw[rel,{Diamond[length=2mm]}-] (BlockDefinition) to node[midway,fill=white,inner sep=2pt,font=\scriptsize,text=BookMuted,align=center] {describes / 1 : 1} (BlockMetadata);
\draw[rel,{Diamond[length=2mm]}-] (SchemaBoundBlock) to node[midway,fill=white,inner sep=2pt,font=\scriptsize,text=BookMuted,align=center] {wraps / 1 : 1} (BlockDefinition);
\draw[rel,{Diamond[length=2mm]}-] (SchemaBoundBlock) to node[midway,fill=white,inner sep=2pt,font=\scriptsize,text=BookMuted,align=center] {validates / 1 : 1} (BlockSchemaAdapter);
\end{tikzpicture}}
\par\smallskip{\small\color{BookMuted}Definition and schema ownership}
\end{center}

\clearpage
\begin{center}
\BookDiagram[.9\linewidth]{%
\begin{tikzpicture}[classbox/.style={draw=BookHeading,fill=BookPanel,line width=.6pt,rounded corners=1.5pt,text width=37mm,align=left,inner sep=5pt,font=\scriptsize},rel/.style={-{Latex[length=2mm]},draw=BookHeading,line width=.6pt}]
\node[classbox] (BlockMetadata) at (0,-7.5) {\textbf{BlockMetadata}\\[-2pt]\rule{\linewidth}{.3pt}\\id\\name\\purpose};
\node[classbox] (BlockSlot) at (7,0) {\textbf{BlockSlot}\\[-2pt]\rule{\linewidth}{.3pt}\\id\\required\\description};
\node[classbox] (BlockVariantMetadata) at (7,-2.5) {\textbf{BlockVariantMetadata}\\[-2pt]\rule{\linewidth}{.3pt}\\id\\name\\purpose};
\node[classbox] (BlockConfigOption) at (7,-5) {\textbf{BlockConfigOption}\\[-2pt]\rule{\linewidth}{.3pt}\\id\\kind\\required};
\node[classbox] (ComposedComponent) at (7,-7.5) {\textbf{ComposedComponent}\\[-2pt]\rule{\linewidth}{.3pt}\\componentId};
\node[classbox] (OutputMode) at (7,-10) {\textbf{OutputMode}\\[-2pt]\rule{\linewidth}{.3pt}\\mode};
\node[classbox] (SemanticFact) at (7,-12.5) {\textbf{SemanticFact}\\[-2pt]\rule{\linewidth}{.3pt}\\kind\\key};
\node[classbox] (StoryCapture) at (7,-15) {\textbf{StoryCapture}\\[-2pt]\rule{\linewidth}{.3pt}\\storyId};
\draw[rel,{Diamond[length=2mm]}-] (BlockMetadata) to node[midway,fill=white,inner sep=2pt,font=\scriptsize,text=BookMuted,align=center] {required slots / 1 : 1..*\\optional slots / 1 : 0..*} (BlockSlot);
\draw[rel,{Diamond[length=2mm]}-] (BlockMetadata) to node[midway,fill=white,inner sep=2pt,font=\scriptsize,text=BookMuted,align=center] {variants / 1 : 0..*} (BlockVariantMetadata);
\draw[rel,{Diamond[length=2mm]}-] (BlockMetadata) to node[midway,fill=white,inner sep=2pt,font=\scriptsize,text=BookMuted,align=center] {config / 1 : 0..*} (BlockConfigOption);
\draw[rel,{Diamond[open,length=2mm]}-] (BlockMetadata) to node[midway,fill=white,inner sep=2pt,font=\scriptsize,text=BookMuted,align=center] {composes / 1 : 0..*} (ComposedComponent);
\draw[rel,{Diamond[length=2mm]}-] (BlockMetadata) to node[midway,fill=white,inner sep=2pt,font=\scriptsize,text=BookMuted,align=center] {lowers to / 1 : 1..*} (OutputMode);
\draw[rel,{Diamond[length=2mm]}-] (BlockMetadata) to node[midway,fill=white,inner sep=2pt,font=\scriptsize,text=BookMuted,align=center] {exposes / 1 : 0..*} (SemanticFact);
\draw[rel,{Diamond[open,length=2mm]}-] (BlockMetadata) to node[midway,fill=white,inner sep=2pt,font=\scriptsize,text=BookMuted,align=center] {proves / 1 : 0..*} (StoryCapture);
\end{tikzpicture}}
\par\smallskip{\small\color{BookMuted}Slots, lowerings and proof metadata}
\end{center}
\clearpage

\clearpage
\begin{center}
\BookDiagram[.9\linewidth]{%
\begin{tikzpicture}[classbox/.style={draw=BookHeading,fill=BookPanel,line width=.6pt,rounded corners=1.5pt,text width=37mm,align=left,inner sep=5pt,font=\scriptsize},rel/.style={-{Latex[length=2mm]},draw=BookHeading,line width=.6pt}]
\node[classbox] (BlockVariantMetadata) at (0,0) {\textbf{BlockVariantMetadata}\\[-2pt]\rule{\linewidth}{.3pt}\\id\\name\\purpose};
\node[classbox] (BlockSlot) at (-3,-3.5) {\textbf{BlockSlot}\\[-2pt]\rule{\linewidth}{.3pt}\\id\\required\\description};
\node[classbox] (SemanticFact) at (3,-3.5) {\textbf{SemanticFact}\\[-2pt]\rule{\linewidth}{.3pt}\\kind\\key};
\draw[rel,{Diamond[length=2mm]}-] (BlockVariantMetadata) to node[midway,fill=white,inner sep=2pt,font=\scriptsize,text=BookMuted,align=center] {variant slots / 1 : 0..*} (BlockSlot);
\draw[rel,{Diamond[length=2mm]}-] (BlockVariantMetadata) to node[midway,fill=white,inner sep=2pt,font=\scriptsize,text=BookMuted,align=center] {variant facts / 1 : 0..*} (SemanticFact);
\end{tikzpicture}}
\par\smallskip{\small\color{BookMuted}Variant-specific slots and facts}
\end{center}
\clearpage

Block package manifests expose definitions. Block metadata owns the contract. Variants, config, modes, semantic facts, and story captures hang off the block instead of being inferred from rendered output.

Metadata sketch:

\begin{CodeBlock}
type BlockScale =
  | 'app'
  | 'workspace'
  | 'section'
  | 'panel'
  | 'control'
  | 'inline'
  | 'item'
  | 'data'
  | 'diagnostic'
  | 'overlay';

interface BlockMetadataDocs {
  readonly summary: string;
  readonly useWhen?: readonly string[];
  readonly avoidWhen?: readonly string[];
  readonly relatedDocs?: readonly string[];
}

interface BlockSlot {
  readonly id: string;
  readonly label?: string;
  readonly required?: boolean;
  readonly description?: string;
}

interface BlockMetadata {
  readonly packageName: string;
  readonly blockName: string;
  readonly family: string;
  readonly scale: BlockScale;
  readonly modes: readonly ('interactive' | 'static' | 'pipe' | 'accessible')[];
  readonly docs: BlockMetadataDocs;
  readonly sourcePath?: string;
  readonly slots: readonly BlockSlot[];
  readonly variants?: readonly BlockVariantMetadata[];
  readonly configOptions?: readonly BlockConfigOption[];
  readonly composedComponents?: readonly string[];
  readonly semanticFacts?: readonly ModeLoweringFact[];
  readonly storyIds?: readonly string[];
  readonly examples?: readonly BlockExample[];
  readonly tags?: readonly string[];
}

interface BlockVariantMetadata {
  readonly id: string;
  readonly label: string;
  readonly description?: string;
  readonly requiredSlots?: readonly string[];
  readonly optionalSlots?: readonly string[];
  readonly facts?: readonly ModeLoweringFact[];
}

interface BlockConfigOption {
  readonly id: string;
  readonly label?: string;
  readonly kind: 'boolean' | 'enum' | 'number' | 'string' | 'adapter';
  readonly required?: boolean;
  readonly description?: string;
  readonly values?: readonly string[];
}

interface BlockExample {
  readonly id: string;
  readonly label: string;
  readonly path?: string;
  readonly command?: string;
}

type DeepReadonly<T> = T extends readonly (infer Item)[]
  ? readonly DeepReadonly<Item>[]
  : T extends object
    ? { readonly [Key in keyof T]: DeepReadonly<T[Key]> }
    : T;

type BlockSchemaResult<Data> =
  | { readonly ok: true; readonly data: DeepReadonly<Data> }
  | { readonly ok: false; readonly issues: readonly BlockSchemaIssue[] };

interface BlockSchemaIssue {
  readonly severity: 'info' | 'warning' | 'error';
  readonly code: string;
  readonly message: string;
  readonly path?: string;
}

interface BlockSchemaDescription {
  readonly requiredFields?: readonly string[];
  readonly optionalFields?: readonly string[];
  readonly redactedFields?: readonly string[];
  readonly facts?: readonly ModeLoweringFact[];
}

interface BlockDefinition<Config = unknown> {
  readonly metadata: BlockMetadata;
  readonly data?: ViewDataContract;
  readonly commands?: readonly CommandIntent[];
  readonly render: (input: BlockRenderInput<Config>) => BlockRenderResult;
}

interface BlockSchemaAdapter<Data> {
  readonly id: string;
  readonly parse: (input: unknown) => BlockSchemaResult<Data>;
  readonly describe?: () => BlockSchemaDescription;
}

interface SchemaBoundBlockDefinition<Data, Config = unknown> {
  readonly block: BlockDefinition<Config>;
  readonly schema: BlockSchemaAdapter<Data>;
  readonly bind: (
    data: DeepReadonly<Data>,
  ) => BlockRenderInput<Config> | {
    readonly input: BlockRenderInput<Config>;
    readonly facts?: readonly ModeLoweringFact[];
  };
}

interface BlockPackageManifest {
  readonly packageName: string;
  readonly version: string;
  readonly bijouPeerRange: string;
  readonly blocks: readonly string[];
  readonly docs?: readonly string[];
  readonly tags?: readonly string[];
}

declare function defineBlock<Config>(
  definition: BlockDefinition<Config>,
): BlockDefinition<Config>;

declare function validateBlockMetadata(
  metadata: BlockMetadata,
): BlockMetadataReport;

declare function blockMetadataSummary(metadata: BlockMetadata): string;

declare function defineSchemaBlock<Data, Config>(
  definition: SchemaBoundBlockDefinition<Data, Config>,
): SchemaBoundBlockDefinition<Data, Config>;

declare function bindSchemaBlockInput<Data, Config>(
  block: SchemaBoundBlockDefinition<Data, Config>,
  input: unknown,
): {
  readonly ok: true;
  readonly input: BlockRenderInput<Config>;
  readonly facts: readonly ModeLoweringFact[];
} | {
  readonly ok: false;
  readonly issues: readonly BlockSchemaIssue[];
  readonly facts: readonly ModeLoweringFact[];
};

declare function defineBlockPackage(
  manifest: BlockPackageManifest,
): BlockPackageManifest;
\end{CodeBlock}

DX-031A lands this as a public metadata-first API in \texttt{@\allowbreak{}f\allowbreak{}l\allowbreak{}y\allowbreak{}i\allowbreak{}n\allowbreak{}g\allowbreak{}r\allowbreak{}o\allowbreak{}b\allowbreak{}o\allowbreak{}t\allowbreak{}s\allowbreak{}/\allowbreak{}b\allowbreak{}i\allowbreak{}j\allowbreak{}o\allowbreak{}u}.

DX-031B Schema-Bound Blocks lands the schema adapter layer as a public boundary contract in \texttt{@\allowbreak{}f\allowbreak{}l\allowbreak{}y\allowbreak{}i\allowbreak{}n\allowbreak{}g\allowbreak{}r\allowbreak{}o\allowbreak{}b\allowbreak{}o\allowbreak{}t\allowbreak{}s\allowbreak{}/\allowbreak{}b\allowbreak{}i\allowbreak{}j\allowbreak{}o\allowbreak{}u}. Schema-bound blocks validate boundary data before producing block render input or binding facts. They do not fetch, subscribe, dispatch, compose runtime views, or render AppShell. The important boundary is that a block and its schema are described before rendering, and their slots, semantic facts, and validation facts are visible to tests, docs, DOGFOOD, and agents.

\section{Public Block API And Distribution}\label{record-128-public-block-api-and-distribution}

Blocks should become a user-facing Bijou API, not an internal DOGFOOD-only catalog.

The public API should let Bijou users and community authors:

\begin{itemize}
\tightlist
\item
  define a block with typed metadata and a render function
\item
  validate block metadata before publishing
\item
  attach examples, story captures, and mode-lowering facts
\item
  package one or more blocks for reuse
\item
  distribute blocks through ordinary NPM packages
\item
  let apps import blocks without hidden registration side effects
\end{itemize}

The first distribution mechanism should be boring: NPM packages with a stable manifest. A future block registry can index NPM metadata, docs, and story captures, but the runtime should not need a central registry to use a block.

Recommended package shape:

\begin{CodeBlock}
@scope/bijou-blocks-commerce
  package.json
  src/index.ts
  blocks/product-list.ts
  blocks/checkout-flow.ts
  stories/product-list.story.ts
  docs/README.md
\end{CodeBlock}

Rules for shared block packages:

\begin{itemize}
\tightlist
\item
  depend on \texttt{@\allowbreak{}f\allowbreak{}l\allowbreak{}y\allowbreak{}i\allowbreak{}n\allowbreak{}g\allowbreak{}r\allowbreak{}o\allowbreak{}b\allowbreak{}o\allowbreak{}t\allowbreak{}s\allowbreak{}/\allowbreak{}b\allowbreak{}i\allowbreak{}j\allowbreak{}o\allowbreak{}u} as a peer dependency
\item
  export block definitions explicitly from the package entrypoint
\item
  avoid global registration or import-time mutation
\item
  include block metadata and package metadata
\item
  include at least one story or fixture per exported block
\item
  declare supported output modes and semantic facts
\item
  keep host effects behind adapters supplied by the consuming app
\item
  never include secret values in metadata, stories, diagnostics, or snapshots
\end{itemize}

Example package entrypoint:

\begin{CodeBlock}
export const productListBlock = defineBlock({
  metadata: productListMetadata,
  render: productList,
});

export const checkoutFlowBlock = defineBlock({
  metadata: checkoutFlowMetadata,
  render: checkoutFlow,
});

export const commerceBlocks = defineBlockPackage({
  packageName: '@example/bijou-blocks-commerce',
  version: '1.0.0',
  bijouPeerRange: '^5.0.0',
  blocks: ['product-list', 'checkout-flow'],
  tags: ['commerce', 'catalog', 'checkout'],
});
\end{CodeBlock}

Runtime-engine posture:

\begin{itemize}
\tightlist
\item
  a block renders into ordinary Bijou visible nodes
\item
  visible nodes still resolve through RE-035 layout envelopes
\item
  interaction still emits commands instead of mutating geometry
\item
  modes lower semantic facts instead of parsing rendered cells
\item
  host adapters own files, credentials, network, storage, and process effects
\end{itemize}

That keeps third-party blocks powerful without making them runtime plugins with unbounded authority.

\section{Relationship To DX-034}\label{record-128-relationship-to-dx-034}

\href{https://github.com/flyingrobots/bijou/blob/cbd70612140e7a450d764b3c7be1c6af878b99bc/docs/design/DX-034-declarative-view-data-binding.md}{DX-034 - Declarative View Data Binding} is the design gate between metadata-only blocks and provider-bound AppShell composition.

DX-031 answers:

\begin{itemize}
\tightlist
\item
  what is a block?
\item
  what slots, variants, modes, facts, stories, and package metadata does it declare?
\item
  how can tooling discover block contracts before rendering?
\end{itemize}

DX-034 answers:

\begin{itemize}
\tightlist
\item
  how does an active block or view declare the data it needs?
\item
  which provider scope satisfies those requirements?
\item
  what immutable snapshot is currently rendered?
\item
  how do user interactions leave the view tree as Commands instead of direct data mutation?
\end{itemize}

Rendered AppShell work should build on both cycles. AppShell needs named block slots, but it also needs a unidirectional data-binding loop so nested navigation, content, inspector, status, and overlay blocks can render provider-backed snapshots without owning business state.

\section{Schema-Bound Blocks}\label{record-128-schema-bound-blocks}

The schema-bound boundary for app builders should be direct:

\begin{enumerate}
\def\labelenumi{\arabic{enumi}.}
\tightlist
\item
  define or import a schema
\item
  bind that schema to a block
\item
  pass a schema-compliant object
\item
  receive immutable block render input and validation facts
\end{enumerate}

In app code, the target ergonomics should feel like:

\begin{CodeBlock}
const productDetailsBaseBlock = defineBlock({
  metadata: productDetailsMetadata,
  render: productDetails,
});

const productDetailsBlock = defineSchemaBlock({
  block: productDetailsBaseBlock,
  schema: zodBlockSchema(productDetailsSchema),
  bind: product => ({
    input: {
      slots: {
        title: product.name,
        summary: product.summary,
        specs: product.specs,
        reviews: product.reviews,
        primaryAction: { id: 'add-to-cart', label: 'Add to cart' },
      },
      config: {
        variant: 'purchase-detail',
      },
    },
    facts: [{ kind: 'entity', key: 'product', value: product.id }],
  }),
});

bindSchemaBlockInput(productDetailsBlock, {
  id: 'bijou-pro',
  name: 'Bijou Pro',
  summary: 'Runtime truth for terminal apps.',
  specs: [{ label: 'Mode', value: 'interactive' }],
  reviews: [],
});
\end{CodeBlock}

That is the real DX target for this layer: give the schema-bound block an unknown object and it should validate the boundary before producing immutable render input, validation issues, and lower-mode facts. Rendering remains the ordinary block renderer's job.

Zod can become the first optional authoring adapter because it is familiar in the TypeScript ecosystem and Bijou already uses it in MCP tooling. The core block contract is schema-adapter shaped so the runtime does not require Zod as a hard dependency.

Recommended split:

\begin{itemize}
\tightlist
\item
  core block API accepts \texttt{B\allowbreak{}l\allowbreak{}o\allowbreak{}c\allowbreak{}k\allowbreak{}S\allowbreak{}c\allowbreak{}h\allowbreak{}e\allowbreak{}m\allowbreak{}a\allowbreak{}A\allowbreak{}d\allowbreak{}a\allowbreak{}p\allowbreak{}t\allowbreak{}e\allowbreak{}r\allowbreak{}<\allowbreak{}D\allowbreak{}a\allowbreak{}t\allowbreak{}a\allowbreak{}>}
\item
  optional Zod helper adapts Zod schemas into that adapter
\item
  future helpers can adapt JSON Schema, Valibot, TypeBox, or generated schemas
\item
  block packages can choose the adapter they publish with
\item
  tooling can inspect schema facts when the adapter exposes them
\end{itemize}

Zod adapter sketch:

\begin{CodeBlock}
type ZodLikeResult<Data> =
  | { readonly success: true; readonly data: Data }
  | { readonly success: false; readonly error: unknown };

interface ZodLikeSchema<Data> {
  readonly safeParse: (input: unknown) => ZodLikeResult<Data>;
}

declare function zodBlockSchema<Data>(
  schema: ZodLikeSchema<Data>,
): BlockSchemaAdapter<Data>;
\end{CodeBlock}

Schema-bound blocks should standardize:

\begin{itemize}
\tightlist
\item
  validation errors into \texttt{E\allowbreak{}r\allowbreak{}r\allowbreak{}o\allowbreak{}r\allowbreak{}S\allowbreak{}c\allowbreak{}r\allowbreak{}e\allowbreak{}e\allowbreak{}n}, \texttt{B\allowbreak{}a\allowbreak{}n\allowbreak{}n\allowbreak{}e\allowbreak{}r}, or field-level facts
\item
  defaults into config and slots
\item
  object fields into named block slots
\item
  object relationships into semantic facts
\item
  mode-lowering facts without parsing rendered output
\item
  docs, examples, and story captures from the same schema-bound fixture data
\end{itemize}

Rules:

\begin{itemize}
\tightlist
\item
  schema validation happens before rendering
\item
  invalid data produces explicit validation facts, not partial mystery UI
\item
  schema-bound blocks validate boundary data and do not fetch, subscribe, dispatch, compose runtime views, or render AppShell
\item
  schema-bound blocks still render through ordinary Bijou layout envelopes
\item
  schemas may describe data shape, but blocks still own presentation choices
\item
  secret fields must be redacted before metadata, diagnostics, or captures
\end{itemize}

\section{Compiler Path: Schema To Blocks}\label{record-128-compiler-path-schema-to-blocks}

The schema-bound API opens a larger tooling path:

% Presentation of original arrow notation. See flow-manifest.json.
\begin{center}
\begin{tikzpicture}[book flow node/.style={draw=BookHeading,line width=.65pt,fill=BookPanel,rounded corners=1.5pt,align=center,inner sep=5pt,minimum height=10mm,font=\fontsize{9}{11}\selectfont,execute at begin node={\hyphenpenalty=10000\exhyphenpenalty=10000}},book flow edge/.style={-{Latex[length=2mm]},draw=BookHeading,line width=.7pt}]
\node[book flow node,text width=96mm] (n0) at (0,0.0) {Graph\allowbreak{}QL schema /\allowbreak{} operations};
\node[book flow node,text width=96mm] (n1) at (0,-1.55) {semantic model};
\node[book flow node,text width=96mm] (n2) at (0,-3.1) {Type\allowbreak{}Script types};
\node[book flow node,text width=96mm] (n3) at (0,-4.65) {Zod or schema adapter};
\node[book flow node,text width=96mm] (n4) at (0,-6.2) {block metadata};
\node[book flow node,text width=96mm] (n5) at (0,-7.75) {schema-bound block package};
\node[book flow node,text width=96mm] (n6) at (0,-9.3) {stories, captures, docs, and tests};
\draw[book flow edge] (n0) -- (n1);
\draw[book flow edge] (n1) -- (n2);
\draw[book flow edge] (n2) -- (n3);
\draw[book flow edge] (n3) -- (n4);
\draw[book flow edge] (n4) -- (n5);
\draw[book flow edge] (n5) -- (n6);
\end{tikzpicture}
\end{center}


This is not first-slice block runtime work, but it is a coherent future lane. \texttt{\textasciitilde{}\allowbreak{}/\allowbreak{}g\allowbreak{}i\allowbreak{}t\allowbreak{}/\allowbreak{}w\allowbreak{}e\allowbreak{}s\allowbreak{}l\allowbreak{}e\allowbreak{}y} already points in this direction: it has GraphQL schema lowering, operation catalogs, TypeScript emission, and dialect seams that mention future SQLite and other SQL targets. Bijou should be able to consume that kind of semantic output and generate useful block packages.

Potential compiler outputs:

\begin{itemize}
\tightlist
\item
  Zod schemas or generic \texttt{B\allowbreak{}l\allowbreak{}o\allowbreak{}c\allowbreak{}k\allowbreak{}S\allowbreak{}c\allowbreak{}h\allowbreak{}e\allowbreak{}m\allowbreak{}a\allowbreak{}A\allowbreak{}d\allowbreak{}a\allowbreak{}p\allowbreak{}t\allowbreak{}e\allowbreak{}r} objects
\item
  TypeScript data types
\item
  default block choices by schema shape
\item
  candidate block metadata
\item
  fixture data and story captures
\item
  mode-lowering semantic facts
\item
  docs tables for fields, relationships, and actions
\end{itemize}

Useful mappings:

\begin{itemize}
\tightlist
\item
  GraphQL object type -\textgreater{} \texttt{P\allowbreak{}r\allowbreak{}o\allowbreak{}d\allowbreak{}u\allowbreak{}c\allowbreak{}t\allowbreak{}D\allowbreak{}e\allowbreak{}t\allowbreak{}a\allowbreak{}i\allowbreak{}l\allowbreak{}s}, \texttt{U\allowbreak{}s\allowbreak{}e\allowbreak{}r\allowbreak{}I\allowbreak{}n\allowbreak{}f\allowbreak{}o}, \texttt{A\allowbreak{}p\allowbreak{}p\allowbreak{}l\allowbreak{}i\allowbreak{}c\allowbreak{}a\allowbreak{}t\allowbreak{}i\allowbreak{}o\allowbreak{}n\allowbreak{}C\allowbreak{}a\allowbreak{}r\allowbreak{}d}
\item
  GraphQL list query -\textgreater{} \texttt{P\allowbreak{}r\allowbreak{}o\allowbreak{}d\allowbreak{}u\allowbreak{}c\allowbreak{}t\allowbreak{}L\allowbreak{}i\allowbreak{}s\allowbreak{}t}, \texttt{G\allowbreak{}r\allowbreak{}i\allowbreak{}d}, \texttt{B\allowbreak{}l\allowbreak{}o\allowbreak{}g\allowbreak{}G\allowbreak{}r\allowbreak{}i\allowbreak{}d}, \texttt{L\allowbreak{}e\allowbreak{}a\allowbreak{}d\allowbreak{}e\allowbreak{}r\allowbreak{}b\allowbreak{}o\allowbreak{}a\allowbreak{}r\allowbreak{}d}
\item
  GraphQL mutation -\textgreater{} \texttt{F\allowbreak{}o\allowbreak{}r\allowbreak{}m\allowbreak{}P\allowbreak{}a\allowbreak{}n\allowbreak{}e\allowbreak{}l}, \texttt{C\allowbreak{}h\allowbreak{}e\allowbreak{}c\allowbreak{}k\allowbreak{}o\allowbreak{}u\allowbreak{}t\allowbreak{}F\allowbreak{}l\allowbreak{}o\allowbreak{}w}, \texttt{B\allowbreak{}u\allowbreak{}t\allowbreak{}t\allowbreak{}o\allowbreak{}n\allowbreak{}S\allowbreak{}e\allowbreak{}t}
\item
  GraphQL enum -\textgreater{} select, segmented control, or \texttt{S\allowbreak{}e\allowbreak{}t\allowbreak{}t\allowbreak{}i\allowbreak{}n\allowbreak{}g\allowbreak{}s\allowbreak{}P\allowbreak{}a\allowbreak{}n\allowbreak{}e\allowbreak{}l} row
\item
  GraphQL relationship -\textgreater{} nested section, \texttt{C\allowbreak{}o\allowbreak{}m\allowbreak{}m\allowbreak{}e\allowbreak{}n\allowbreak{}t\allowbreak{}s\allowbreak{}T\allowbreak{}h\allowbreak{}r\allowbreak{}e\allowbreak{}a\allowbreak{}d}, or detail pane
\item
  database table -\textgreater{} grid/list/detail blocks
\item
  database view -\textgreater{} read-only dashboard or data visualization block
\end{itemize}

Local storage can enter through the same seam. A SQLite-backed app, local JSON store, file-backed cache, or generated persistence adapter can provide the data object. The block should not care whether the object came from GraphQL, SQLite, REST, a file, or memory after it passes the schema adapter.

The compiler should produce suggestions, not magic authority. Generated blocks still need:

\begin{itemize}
\tightlist
\item
  explicit metadata
\item
  explicit schema adapters
\item
  explicit slot binding
\item
  explicit mode-lowering facts
\item
  story captures with representative data
\item
  human-reviewable output
\end{itemize}

The win is speed. A maintainer should be able to point a compiler at a semantic schema and get a boring, inspectable first draft of a Bijou app: lists, details, forms, dashboards, charts, empty/error states, and docs surfaces that already obey the block contract.

\section{Tooling Payoff}\label{record-128-tooling-payoff}

A metadata-first block API should make future Bijou tooling much cheaper.

Tools should be able to index block packages and answer:

\begin{itemize}
\tightlist
\item
  which blocks exist?
\item
  which slots are required?
\item
  which variants are supported?
\item
  which modes does a block lower to?
\item
  which semantic facts does a block promise?
\item
  which stories or captures prove the block?
\item
  which components, themes, or adapters does the block depend on?
\end{itemize}

That directly supports:

\begin{itemize}
\tightlist
\item
  a Storybook-like block and component workstation
\item
  theme token inspectors
\item
  layout inspectors
\item
  focus and input routing inspectors
\item
  mode-lowering reports
\item
  DOGFOOD block galleries
\item
  MCP docs and block search payloads
\item
  fixture-to-docs promotion tooling
\item
  block package compatibility reports
\item
  visual and semantic regression review
\end{itemize}

The important constraint is that tooling reads declared block facts first. It may render stories, capture screenshots, inspect surfaces, or run assertions afterward, but it should not need to scrape terminal output just to discover what a block is.

\section{Variants And Configuration}\label{record-128-variants-and-configuration}

Blocks should offer named variants when the job is stable but the surface shape or input contract changes.

A variant is still the same block. It should keep the same broad purpose while declaring any changed slots, focus rules, semantic facts, and mode-lowering behavior.

Configuration is different from variants. Configuration changes parameters inside a chosen variant.

Examples:

\begin{itemize}
\tightlist
\item
  \texttt{A\allowbreak{}p\allowbreak{}p\allowbreak{}S\allowbreak{}h\allowbreak{}e\allowbreak{}l\allowbreak{}l} can offer docs, workbench, and focused-task variants.
\item
  \texttt{C\allowbreak{}o\allowbreak{}m\allowbreak{}m\allowbreak{}a\allowbreak{}n\allowbreak{}d\allowbreak{}P\allowbreak{}a\allowbreak{}l\allowbreak{}e\allowbreak{}t\allowbreak{}t\allowbreak{}e} can offer commands-only, commands-with-preview, and destination-search variants.
\item
  \texttt{H\allowbreak{}e\allowbreak{}r\allowbreak{}o} can offer compact, branded, docs-landing, and playable-demo variants.
\item
  \texttt{L\allowbreak{}o\allowbreak{}a\allowbreak{}d\allowbreak{}i\allowbreak{}n\allowbreak{}g\allowbreak{}S\allowbreak{}c\allowbreak{}r\allowbreak{}e\allowbreak{}e\allowbreak{}n} can offer indeterminate, determinate, and staged-bootstrap variants.
\item
  \texttt{T\allowbreak{}i\allowbreak{}t\allowbreak{}l\allowbreak{}e\allowbreak{}S\allowbreak{}c\allowbreak{}r\allowbreak{}e\allowbreak{}e\allowbreak{}n} can offer press-enter, loading-gate, and auth-gateway variants.
\item
  \texttt{P\allowbreak{}e\allowbreak{}r\allowbreak{}f\allowbreak{}H\allowbreak{}u\allowbreak{}d} can offer compact overlay, detailed panel, and recorder variants.
\item
  \texttt{L\allowbreak{}o\allowbreak{}g\allowbreak{}i\allowbreak{}n\allowbreak{}F\allowbreak{}o\allowbreak{}r\allowbreak{}m} can offer manual-password, one-time-code, and credential-provider variants.
\item
  \texttt{S\allowbreak{}i\allowbreak{}g\allowbreak{}n\allowbreak{}U\allowbreak{}p\allowbreak{}F\allowbreak{}o\allowbreak{}r\allowbreak{}m} can offer password, invite-code, and provider-handoff variants.
\end{itemize}

Rules:

\begin{itemize}
\tightlist
\item
  each variant must declare required slots and semantic facts
\item
  each variant must lower across interactive, static, pipe, and accessible modes
\item
  config must be inspectable without exposing secret values
\item
  unsupported variants should fail clearly instead of silently degrading
\item
  host adapters provide credentials, files, network access, and other effects
\end{itemize}

Secret-bearing config is not block metadata. The metadata may name a credential provider capability, but the secret value stays in host-owned state.

Auth blocks should use a provider seam rather than hard-code a vault package. A host can adapt a provider such as \texttt{@\allowbreak{}g\allowbreak{}i\allowbreak{}t\allowbreak{}-\allowbreak{}s\allowbreak{}t\allowbreak{}u\allowbreak{}n\allowbreak{}t\allowbreak{}s\allowbreak{}/\allowbreak{}v\allowbreak{}a\allowbreak{}u\allowbreak{}l\allowbreak{}t} behind that seam, but the block contract should only see a named credential provider and its safe status facts.

Sketch:

\begin{CodeBlock}
type LoginCredentialMode =
  | 'manual-password'
  | 'one-time-code'
  | 'credential-provider';

interface LoginBlockConfig {
  readonly credentialMode: LoginCredentialMode;
  readonly credentialProviderId?: string;
  readonly allowRememberIdentity?: boolean;
  readonly requireMfaCode?: boolean;
}
\end{CodeBlock}

\section{Existing Proto-Blocks}\label{record-128-existing-proto-blocks}

Some Bijou componentry is already block-shaped. DX-031 should inventory and promote that work instead of rebuilding it under new names.

A component or component group is a proto-block when it already has:

\begin{itemize}
\tightlist
\item
  a recognizable product job
\item
  multiple composed lower-level pieces
\item
  stable visible regions
\item
  reusable input behavior
\item
  mode-lowering expectations
\item
  diagnostic or semantic facts worth testing
\end{itemize}

The promotion path should be:

\begin{enumerate}
\def\labelenumi{\arabic{enumi}.}
\tightlist
\item
  identify the existing component or story
\item
  name the block job it already performs
\item
  declare slots, variants, config, modes, and semantic facts
\item
  add story captures and mode-lowering assertions
\item
  keep the implementation if it already fits the contract
\end{enumerate}

The important rule is that ``block'' is a standardization level, not a rewrite mandate. If existing componentry already satisfies the contract, promotion should document and test it first.

\section{Block Scale}\label{record-128-block-scale}

Blocks can exist at different scales. DX-031 should name the scale instead of pretending every useful pattern is a full application shell.

Block scales:

\begin{itemize}
\tightlist
\item
  app blocks: whole app states or persistent shells
\item
  workspace blocks: full app workspaces under a shell
\item
  section blocks: large content or product sections inside an app
\item
  panel blocks: reusable panes, drawers, modals, and diagnostics
\item
  control blocks: repeated input/action compositions
\item
  inline blocks: local status, hints, and compact feedback beside content
\item
  item blocks: repeated cards, comments, rows, and list entries
\item
  data blocks: tabular, chart, or evidence-focused semantic surfaces
\item
  diagnostic blocks: inspectable explanation, health, and instrumentation surfaces
\item
  overlay blocks: transient shell-level feedback or affordances
\end{itemize}

Some names are component-like on purpose. \texttt{B\allowbreak{}u\allowbreak{}t\allowbreak{}t\allowbreak{}o\allowbreak{}n} and \texttt{S\allowbreak{}l\allowbreak{}i\allowbreak{}d\allowbreak{}e\allowbreak{}r} may remain primitive components, while \texttt{B\allowbreak{}u\allowbreak{}t\allowbreak{}t\allowbreak{}o\allowbreak{}n\allowbreak{}S\allowbreak{}e\allowbreak{}t}, \texttt{A\allowbreak{}c\allowbreak{}t\allowbreak{}i\allowbreak{}o\allowbreak{}n\allowbreak{}F\allowbreak{}o\allowbreak{}o\allowbreak{}t\allowbreak{}e\allowbreak{}r}, or \texttt{S\allowbreak{}l\allowbreak{}i\allowbreak{}d\allowbreak{}e\allowbreak{}r\allowbreak{}C\allowbreak{}o\allowbreak{}n\allowbreak{}t\allowbreak{}r\allowbreak{}o\allowbreak{}l} can be blocks when they standardize layout, labels, validation, focus, and mode-lowering facts around those primitives.

\section{Block Families}\label{record-128-block-families}

The first catalog should group blocks by job rather than by implementation package. That keeps the list usable for humans and selectable by agents.

App structure:

\begin{itemize}
\tightlist
\item
  AppShell
\item
  DashboardShell
\item
  HeaderBar
\item
  FooterBar
\item
  NavigationBar
\item
  ReaderSurface
\item
  NavigationList
\item
  TableOfContents
\item
  InspectorPanel
\item
  TitleScreen
\item
  HeroHeader
\item
  Hero
\item
  LoadingScreen
\end{itemize}

Content and marketing sections:

\begin{itemize}
\tightlist
\item
  FeatureBlock
\item
  FeaturesSection
\item
  FAQSection
\item
  Banner
\item
  CallToAction
\item
  PricingTiers
\item
  StatsPanel
\item
  ComparisonPanel
\item
  ErrorScreen
\item
  CommentsThread
\item
  Carousel
\item
  Timeline
\end{itemize}

Command and assistance:

\begin{itemize}
\tightlist
\item
  CommandPalette
\item
  HelpPanel
\item
  SettingsPanel
\end{itemize}

Transient and blocking overlays:

\begin{itemize}
\tightlist
\item
  NotificationStack
\item
  ModalDialogStack
\item
  QuitConfirmationDialog
\end{itemize}

Forms and identity:

\begin{itemize}
\tightlist
\item
  UserInfo
\item
  FormPanel
\item
  LoginForm
\item
  SignUpForm
\end{itemize}

Controls:

\begin{itemize}
\tightlist
\item
  ButtonSet
\item
  SliderControl
\end{itemize}

Data and workbench candidates:

\begin{itemize}
\tightlist
\item
  Grid
\item
  WorkbenchTable
\item
  FormFlow
\item
  ActivityLog
\item
  EmptyErrorState
\item
  DetailPane
\item
  SearchResultsPane
\item
  FileExplorerPane
\item
  ApplicationCard
\item
  ApplicationCardGrid
\item
  DataVisualizationPanel
\item
  ProductList
\item
  ProductDetails
\item
  ProductCategory
\item
  ProductSpecsSheet
\item
  BlogGrid
\item
  BlogPost
\item
  Leaderboard
\item
  TerminalPanel
\end{itemize}

Commerce:

\begin{itemize}
\tightlist
\item
  ShoppingCart
\item
  CheckoutFlow
\item
  PaymentMethods
\item
  OfferModal
\item
  ReviewsBlock
\end{itemize}

Communication:

\begin{itemize}
\tightlist
\item
  ChatPanel
\end{itemize}

Runtime diagnostics:

\begin{itemize}
\tightlist
\item
  PerfHud
\item
  LayoutInspectorPanel
\item
  InputRoutingInspector
\item
  SurfaceDiffPanel
\end{itemize}

\section{Initial Block Catalog}\label{record-128-initial-block-catalog}

\subsection{AppShell}\label{record-128-appshell}

Standard top-level app frame.

Variants:

\begin{itemize}
\tightlist
\item
  docs shell
\item
  workbench shell
\item
  focused-task shell
\end{itemize}

Slots:

\begin{itemize}
\tightlist
\item
  header
\item
  navigation
\item
  content
\item
  inspector
\item
  status
\item
  overlays
\end{itemize}

Behavior:

\begin{itemize}
\tightlist
\item
  owns the top-level app layout
\item
  switches between narrow, standard, and wide variants
\item
  keeps content primary in constrained terminals
\item
  exposes named regions for focus and accessibility
\end{itemize}

\subsection{ReaderSurface}\label{record-128-readersurface}

Standard long-form reading surface for docs, guides, release notes, and prose.

Variants:

\begin{itemize}
\tightlist
\item
  article reader
\item
  reference page
\item
  release notes
\end{itemize}

Slots:

\begin{itemize}
\tightlist
\item
  title
\item
  breadcrumbs
\item
  body
\item
  table of contents
\item
  source or provenance metadata
\item
  actions
\end{itemize}

Behavior:

\begin{itemize}
\tightlist
\item
  preserves full semantic content in pipe mode
\item
  reports current section and scroll facts in accessible mode
\item
  can live inside \texttt{A\allowbreak{}p\allowbreak{}p\allowbreak{}S\allowbreak{}h\allowbreak{}e\allowbreak{}l\allowbreak{}l} as the primary content region
\end{itemize}

\subsection{NavigationList}\label{record-128-navigationlist}

Grouped navigation list for sidebars, drawers, and quick navigation surfaces.

Variants:

\begin{itemize}
\tightlist
\item
  grouped navigation
\item
  flat navigation
\item
  filtered navigation
\end{itemize}

Slots:

\begin{itemize}
\tightlist
\item
  groups
\item
  items
\item
  active item
\item
  badges
\item
  filter
\item
  empty state
\end{itemize}

Behavior:

\begin{itemize}
\tightlist
\item
  keeps the active item visible when possible
\item
  exposes grouped navigation facts
\item
  lowers to source-order destinations in pipe mode
\end{itemize}

\subsection{InspectorPanel}\label{record-128-inspectorpanel}

Structured facts and contextual actions for selected entities, diagnostics, layout nodes, focus maps, runtime records, or theme reports.

Variants:

\begin{itemize}
\tightlist
\item
  facts panel
\item
  diagnostics panel
\item
  action panel
\end{itemize}

Slots:

\begin{itemize}
\tightlist
\item
  title
\item
  summary
\item
  sections
\item
  facts
\item
  warnings
\item
  actions
\item
  footer
\end{itemize}

Behavior:

\begin{itemize}
\tightlist
\item
  renders stable key-value facts
\item
  exposes warnings and actions semantically
\item
  works as a pane, drawer, or modal body
\end{itemize}

\subsection{PerfHud}\label{record-128-perfhud}

Performance and runtime health surface for development builds, diagnostics, benchmarking, and DOGFOOD instrumentation.

Variants:

\begin{itemize}
\tightlist
\item
  compact overlay
\item
  detailed panel
\item
  recorder
\item
  budget warning strip
\end{itemize}

Config:

\begin{itemize}
\tightlist
\item
  metric set
\item
  sample interval
\item
  budget thresholds
\item
  history depth
\item
  overlay position
\end{itemize}

Slots:

\begin{itemize}
\tightlist
\item
  title
\item
  frame metrics
\item
  render metrics
\item
  memory metrics
\item
  event-loop metrics
\item
  budget warnings
\item
  history sparkline
\item
  actions
\end{itemize}

Behavior:

\begin{itemize}
\tightlist
\item
  can be promoted from existing Perf HUD componentry
\item
  renders as overlay, panel, drawer, or status strip
\item
  exposes current metrics as semantic facts
\item
  lowers to stable metric lines in static and pipe modes
\item
  reports threshold state in accessible mode
\end{itemize}

\subsection{CommandPalette}\label{record-128-commandpalette}

Searchable command and destination picker.

Variants:

\begin{itemize}
\tightlist
\item
  commands only
\item
  commands with preview
\item
  destination search
\end{itemize}

Slots:

\begin{itemize}
\tightlist
\item
  query
\item
  grouped results
\item
  selected result
\item
  preview
\item
  footer hints
\item
  empty state
\end{itemize}

Behavior:

\begin{itemize}
\tightlist
\item
  keyboard-first command execution
\item
  optional preview in wide layouts
\item
  accessible combobox/listbox-style facts
\item
  pipe mode can emit matching command ids and titles
\end{itemize}

\subsection{HelpPanel}\label{record-128-helppanel}

Standard help and keybinding surface.

Variants:

\begin{itemize}
\tightlist
\item
  global help
\item
  contextual help
\item
  searchable keybindings
\end{itemize}

Slots:

\begin{itemize}
\tightlist
\item
  title
\item
  context summary
\item
  command groups
\item
  shortcuts
\item
  search/filter
\item
  footer hints
\end{itemize}

Behavior:

\begin{itemize}
\tightlist
\item
  can render as a panel, drawer, or modal
\item
  groups commands by context
\item
  exposes available actions in accessible mode
\end{itemize}

\subsection{SettingsPanel}\label{record-128-settingspanel}

Standard preferences surface for theme, shell, density, mode, and app settings.

Variants:

\begin{itemize}
\tightlist
\item
  compact settings
\item
  sectioned settings
\item
  full-screen narrow settings
\end{itemize}

Slots:

\begin{itemize}
\tightlist
\item
  sections
\item
  rows
\item
  active section
\item
  descriptions
\item
  actions
\end{itemize}

Rows:

\begin{itemize}
\tightlist
\item
  toggle
\item
  select
\item
  segmented choice
\item
  theme swatch
\item
  numeric stepper
\item
  readonly fact
\item
  danger action
\end{itemize}

Behavior:

\begin{itemize}
\tightlist
\item
  aligns values consistently
\item
  supports drawer or full-screen narrow presentation
\item
  lowers to \texttt{k\allowbreak{}e\allowbreak{}y\allowbreak{}=\allowbreak{}v\allowbreak{}a\allowbreak{}l\allowbreak{}u\allowbreak{}e} facts in pipe mode
\end{itemize}

\subsection{NotificationStack}\label{record-128-notificationstack}

Stacked notifications and transient status messages.

Variants:

\begin{itemize}
\tightlist
\item
  transient toasts
\item
  persistent history
\item
  mixed severity stack
\end{itemize}

Slots:

\begin{itemize}
\tightlist
\item
  notifications
\item
  active toast
\item
  history
\item
  actions
\end{itemize}

Behavior:

\begin{itemize}
\tightlist
\item
  keeps transient display separate from notification history
\item
  exposes severity, title, body, and action facts
\item
  should eventually route hit-testing through retained layout
\end{itemize}

\subsection{ModalDialogStack}\label{record-128-modaldialogstack}

Standard stacked modal surface for blocking overlays.

Variants:

\begin{itemize}
\tightlist
\item
  single blocking modal
\item
  nested modal stack
\item
  confirmation flow
\end{itemize}

Slots:

\begin{itemize}
\tightlist
\item
  modal stack
\item
  active modal
\item
  backdrop
\item
  title
\item
  body
\item
  actions
\end{itemize}

Behavior:

\begin{itemize}
\tightlist
\item
  topmost modal receives input first
\item
  lower app surfaces are blocked while active
\item
  focus is trapped inside the active modal
\item
  accessible mode reports modal title, body, and available actions
\end{itemize}

\subsection{QuitConfirmationDialog}\label{record-128-quitconfirmationdialog}

Specialized modal for quitting, closing, or abandoning dirty work.

Variants:

\begin{itemize}
\tightlist
\item
  quit cleanly
\item
  quit with dirty work
\item
  save before quit
\end{itemize}

Slots:

\begin{itemize}
\tightlist
\item
  title
\item
  message
\item
  dirty summary
\item
  confirm action
\item
  cancel action
\item
  optional save action
\end{itemize}

Behavior:

\begin{itemize}
\tightlist
\item
  defaults to cancel as the safer action
\item
  reports destructive intent clearly
\item
  can be hosted by \texttt{M\allowbreak{}o\allowbreak{}d\allowbreak{}a\allowbreak{}l\allowbreak{}D\allowbreak{}i\allowbreak{}a\allowbreak{}l\allowbreak{}o\allowbreak{}g\allowbreak{}S\allowbreak{}t\allowbreak{}a\allowbreak{}c\allowbreak{}k}
\end{itemize}

\subsection{LoadingScreen}\label{record-128-loadingscreen}

Readiness surface for startup, connection, indexing, generation, or other bounded and unbounded waits.

Variants:

\begin{itemize}
\tightlist
\item
  indeterminate loading
\item
  determinate progress
\item
  staged bootstrap
\item
  retryable failure
\end{itemize}

Config:

\begin{itemize}
\tightlist
\item
  progress mode
\item
  stage labels
\item
  cancel policy
\item
  retry policy
\item
  timeout policy
\end{itemize}

Slots:

\begin{itemize}
\tightlist
\item
  title
\item
  current stage
\item
  progress value
\item
  progress label
\item
  detail message
\item
  cancel action
\item
  retry action
\end{itemize}

Behavior:

\begin{itemize}
\tightlist
\item
  exposes readiness and progress facts without relying on animation
\item
  lowers to stable stage/progress text in static and pipe modes
\item
  reports retryable failure state accessibly
\item
  can be embedded inside \texttt{T\allowbreak{}i\allowbreak{}t\allowbreak{}l\allowbreak{}e\allowbreak{}S\allowbreak{}c\allowbreak{}r\allowbreak{}e\allowbreak{}e\allowbreak{}n}
\end{itemize}

\subsection{TitleScreen}\label{record-128-titlescreen}

App entry surface before the main application shell becomes active.

Variants:

\begin{itemize}
\tightlist
\item
  press enter
\item
  loading gate
\item
  auth gateway
\item
  profile chooser
\end{itemize}

Config:

\begin{itemize}
\tightlist
\item
  entry action
\item
  startup state
\item
  auth requirement
\item
  default next route
\item
  skip policy
\end{itemize}

Slots:

\begin{itemize}
\tightlist
\item
  title
\item
  subtitle
\item
  visual surface
\item
  loading surface
\item
  login form
\item
  sign-up form
\item
  primary action
\item
  secondary actions
\item
  status or version facts
\end{itemize}

Behavior:

\begin{itemize}
\tightlist
\item
  owns the app entry state before \texttt{A\allowbreak{}p\allowbreak{}p\allowbreak{}S\allowbreak{}h\allowbreak{}e\allowbreak{}l\allowbreak{}l} becomes active
\item
  can compose \texttt{H\allowbreak{}e\allowbreak{}r\allowbreak{}o}, \texttt{L\allowbreak{}o\allowbreak{}a\allowbreak{}d\allowbreak{}i\allowbreak{}n\allowbreak{}g\allowbreak{}S\allowbreak{}c\allowbreak{}r\allowbreak{}e\allowbreak{}e\allowbreak{}n}, \texttt{L\allowbreak{}o\allowbreak{}g\allowbreak{}i\allowbreak{}n\allowbreak{}F\allowbreak{}o\allowbreak{}r\allowbreak{}m}, and \texttt{S\allowbreak{}i\allowbreak{}g\allowbreak{}n\allowbreak{}U\allowbreak{}p\allowbreak{}F\allowbreak{}o\allowbreak{}r\allowbreak{}m}
\item
  exposes whether the app is waiting, ready, blocked, or authenticated
\item
  preserves title, readiness, and available actions in lower modes
\end{itemize}

\subsection{Hero}\label{record-128-hero}

Opening/title surface for branded apps, demos, docs landing screens, and playable examples.

Variants:

\begin{itemize}
\tightlist
\item
  compact hero
\item
  branded hero
\item
  docs landing hero
\item
  playable demo hero
\end{itemize}

Slots:

\begin{itemize}
\tightlist
\item
  title
\item
  subtitle
\item
  visual surface
\item
  primary action
\item
  secondary actions
\item
  status or version facts
\end{itemize}

Behavior:

\begin{itemize}
\tightlist
\item
  remains useful when visual ornamentation lowers away
\item
  can host shader or canvas visuals inside a layout-owned envelope
\item
  pipe and accessible modes preserve title, purpose, and actions
\end{itemize}

\subsection{LoginForm}\label{record-128-loginform}

Standard authentication form for email/password, token, or device-code flows.

Variants:

\begin{itemize}
\tightlist
\item
  manual password
\item
  one-time code
\item
  credential provider
\end{itemize}

Config:

\begin{itemize}
\tightlist
\item
  credential mode
\item
  credential provider id
\item
  remember identity policy
\item
  MFA requirement
\end{itemize}

Slots:

\begin{itemize}
\tightlist
\item
  title
\item
  fields
\item
  credential input
\item
  credential provider status
\item
  validation summary
\item
  submit action
\item
  secondary actions
\item
  provider options
\end{itemize}

Behavior:

\begin{itemize}
\tightlist
\item
  never hides validation facts in styling alone
\item
  supports non-interactive pipe/accessibility summaries without exposing secrets
\item
  can request credentials from a host-owned provider seam
\item
  can adapt \texttt{@\allowbreak{}g\allowbreak{}i\allowbreak{}t\allowbreak{}-\allowbreak{}s\allowbreak{}t\allowbreak{}u\allowbreak{}n\allowbreak{}t\allowbreak{}s\allowbreak{}/\allowbreak{}v\allowbreak{}a\allowbreak{}u\allowbreak{}l\allowbreak{}t} through that seam without hard-coding it
\item
  can compose with \texttt{F\allowbreak{}o\allowbreak{}r\allowbreak{}m\allowbreak{}F\allowbreak{}l\allowbreak{}o\allowbreak{}w} later
\end{itemize}

\subsection{SignUpForm}\label{record-128-signupform}

Standard account creation form.

Variants:

\begin{itemize}
\tightlist
\item
  password signup
\item
  invite-code signup
\item
  provider handoff
\end{itemize}

Config:

\begin{itemize}
\tightlist
\item
  required field set
\item
  consent policy
\item
  provider id
\item
  validation mode
\end{itemize}

Slots:

\begin{itemize}
\tightlist
\item
  title
\item
  fields
\item
  validation summary
\item
  terms acknowledgement
\item
  submit action
\item
  secondary actions
\end{itemize}

Behavior:

\begin{itemize}
\tightlist
\item
  separates validation, consent, and submit states
\item
  lowers field requirements and validation facts accessibly
\item
  can share form-row contracts with \texttt{L\allowbreak{}o\allowbreak{}g\allowbreak{}i\allowbreak{}n\allowbreak{}F\allowbreak{}o\allowbreak{}r\allowbreak{}m}
\end{itemize}

\section{Additional Candidate Blocks}\label{record-128-additional-candidate-blocks}

These blocks are useful enough to name now, but should follow the first implementation set unless an existing proto-block is ready to promote.

\subsection{HeaderBar}\label{record-128-headerbar}

Top chrome for application, section, or panel identity.

Variants:

\begin{itemize}
\tightlist
\item
  app header
\item
  section header
\item
  compact header
\item
  hero header
\end{itemize}

Slots:

\begin{itemize}
\tightlist
\item
  title
\item
  subtitle
\item
  breadcrumbs
\item
  primary action
\item
  secondary actions
\item
  status facts
\end{itemize}

Behavior:

\begin{itemize}
\tightlist
\item
  exposes heading level and current location semantically
\item
  can host search, command, help, or account affordances
\item
  lowers to title, location, status, and actions in pipe mode
\end{itemize}

\subsection{FooterBar}\label{record-128-footerbar}

Bottom chrome for status, actions, navigation, or provenance.

Variants:

\begin{itemize}
\tightlist
\item
  status footer
\item
  action footer
\item
  pagination footer
\item
  provenance footer
\end{itemize}

Slots:

\begin{itemize}
\tightlist
\item
  status
\item
  primary action
\item
  secondary actions
\item
  pagination
\item
  provenance
\item
  shortcut hints
\end{itemize}

Behavior:

\begin{itemize}
\tightlist
\item
  keeps status and action facts separate
\item
  can host \texttt{B\allowbreak{}u\allowbreak{}t\allowbreak{}t\allowbreak{}o\allowbreak{}n\allowbreak{}S\allowbreak{}e\allowbreak{}t} without hiding actions in decoration
\item
  lowers shortcut hints only when useful for the output mode
\end{itemize}

\subsection{TableOfContents}\label{record-128-tableofcontents}

Structured navigation for long-form content.

Variants:

\begin{itemize}
\tightlist
\item
  flat contents
\item
  nested contents
\item
  scroll-synced contents
\item
  filtered contents
\end{itemize}

Slots:

\begin{itemize}
\tightlist
\item
  sections
\item
  active section
\item
  progress facts
\item
  filter
\end{itemize}

Behavior:

\begin{itemize}
\tightlist
\item
  exposes document structure and active section
\item
  can live in \texttt{R\allowbreak{}e\allowbreak{}a\allowbreak{}d\allowbreak{}e\allowbreak{}r\allowbreak{}S\allowbreak{}u\allowbreak{}r\allowbreak{}f\allowbreak{}a\allowbreak{}c\allowbreak{}e}, drawer, or narrow navigation mode
\item
  lowers to source-order section links in pipe mode
\end{itemize}

\subsection{HeroHeader}\label{record-128-heroheader}

Header-scale opening region for pages that need title, pitch, status, and actions without becoming a full \texttt{H\allowbreak{}e\allowbreak{}r\allowbreak{}o}.

Variants:

\begin{itemize}
\tightlist
\item
  compact page hero
\item
  product hero header
\item
  docs hero header
\item
  object hero header
\end{itemize}

Slots:

\begin{itemize}
\tightlist
\item
  title
\item
  summary
\item
  media or visual surface
\item
  primary action
\item
  secondary actions
\item
  status facts
\end{itemize}

Behavior:

\begin{itemize}
\tightlist
\item
  keeps the title as the first semantic signal
\item
  can degrade to an ordinary \texttt{H\allowbreak{}e\allowbreak{}a\allowbreak{}d\allowbreak{}e\allowbreak{}r\allowbreak{}B\allowbreak{}a\allowbreak{}r}
\item
  preserves actions and status in lower modes
\end{itemize}

\subsection{FeaturesSection}\label{record-128-featuressection}

Grouped feature or capability presentation.

Variants:

\begin{itemize}
\tightlist
\item
  feature list
\item
  feature grid
\item
  comparison features
\item
  checklist features
\end{itemize}

Slots:

\begin{itemize}
\tightlist
\item
  heading
\item
  feature items
\item
  icons or markers
\item
  supporting text
\item
  actions
\end{itemize}

Behavior:

\begin{itemize}
\tightlist
\item
  exposes each feature as a semantic item
\item
  avoids relying on equal-height visual cards for meaning
\item
  lowers to a stable ordered list
\end{itemize}

\subsection{FAQSection}\label{record-128-faqsection}

Question and answer disclosure surface.

Variants:

\begin{itemize}
\tightlist
\item
  static FAQ
\item
  accordion FAQ
\item
  searchable FAQ
\item
  categorized FAQ
\end{itemize}

Slots:

\begin{itemize}
\tightlist
\item
  questions
\item
  answers
\item
  active question
\item
  categories
\item
  search/filter
\end{itemize}

Behavior:

\begin{itemize}
\tightlist
\item
  exposes expanded/collapsed state
\item
  preserves all answers in pipe mode by policy
\item
  can lower to headings and paragraphs without visual chrome
\end{itemize}

\subsection{Banner}\label{record-128-banner}

Inline announcement, warning, success, or callout surface.

Variants:

\begin{itemize}
\tightlist
\item
  info banner
\item
  warning banner
\item
  success banner
\item
  error banner
\item
  action banner
\end{itemize}

Slots:

\begin{itemize}
\tightlist
\item
  severity
\item
  title
\item
  body
\item
  icon or marker
\item
  actions
\item
  dismiss action
\end{itemize}

Behavior:

\begin{itemize}
\tightlist
\item
  exposes severity as a semantic fact
\item
  does not rely on color alone
\item
  can promote to notification or modal behavior when configured
\end{itemize}

\subsection{ErrorScreen}\label{record-128-errorscreen}

Full-screen or panel-sized error recovery surface.

Variants:

\begin{itemize}
\tightlist
\item
  not found
\item
  permission denied
\item
  crashed
\item
  offline
\item
  retryable failure
\end{itemize}

Slots:

\begin{itemize}
\tightlist
\item
  title
\item
  message
\item
  diagnostic facts
\item
  retry action
\item
  fallback action
\item
  report action
\end{itemize}

Behavior:

\begin{itemize}
\tightlist
\item
  separates user-facing message from diagnostic facts
\item
  exposes recovery actions clearly
\item
  lowers to stable error text and available actions
\end{itemize}

\subsection{UserInfo}\label{record-128-userinfo}

Account, identity, profile, or actor summary.

Variants:

\begin{itemize}
\tightlist
\item
  compact identity
\item
  profile summary
\item
  account switcher item
\item
  permission summary
\end{itemize}

Slots:

\begin{itemize}
\tightlist
\item
  display name
\item
  handle or email
\item
  avatar or marker
\item
  role
\item
  status
\item
  actions
\end{itemize}

Behavior:

\begin{itemize}
\tightlist
\item
  keeps identity facts separate from authentication secrets
\item
  can compose into headers, title screens, settings, and comments
\item
  lowers to name, role, status, and actions
\end{itemize}

\subsection{FormPanel}\label{record-128-formpanel}

General form composition for fields, validation, and actions.

Variants:

\begin{itemize}
\tightlist
\item
  single step form
\item
  multi-step form
\item
  review form
\item
  readonly form summary
\end{itemize}

Slots:

\begin{itemize}
\tightlist
\item
  title
\item
  field groups
\item
  validation summary
\item
  dirty state
\item
  primary action
\item
  secondary actions
\end{itemize}

Behavior:

\begin{itemize}
\tightlist
\item
  exposes required fields, validation state, and dirty state
\item
  can compose \texttt{B\allowbreak{}u\allowbreak{}t\allowbreak{}t\allowbreak{}o\allowbreak{}n\allowbreak{}S\allowbreak{}e\allowbreak{}t}, \texttt{S\allowbreak{}l\allowbreak{}i\allowbreak{}d\allowbreak{}e\allowbreak{}r\allowbreak{}C\allowbreak{}o\allowbreak{}n\allowbreak{}t\allowbreak{}r\allowbreak{}o\allowbreak{}l}, \texttt{L\allowbreak{}o\allowbreak{}g\allowbreak{}i\allowbreak{}n\allowbreak{}F\allowbreak{}o\allowbreak{}r\allowbreak{}m}, and \texttt{S\allowbreak{}i\allowbreak{}g\allowbreak{}n\allowbreak{}U\allowbreak{}p\allowbreak{}F\allowbreak{}o\allowbreak{}r\allowbreak{}m}
\item
  lowers to field facts and validation facts without secret values
\end{itemize}

\subsection{ButtonSet}\label{record-128-buttonset}

Grouped action controls with explicit priority and focus order.

Variants:

\begin{itemize}
\tightlist
\item
  primary and secondary
\item
  toolbar buttons
\item
  segmented action group
\item
  destructive confirmation actions
\end{itemize}

Slots:

\begin{itemize}
\tightlist
\item
  primary action
\item
  secondary actions
\item
  destructive action
\item
  disabled reasons
\item
  shortcut hints
\end{itemize}

Behavior:

\begin{itemize}
\tightlist
\item
  exposes action priority and disabled reasons
\item
  keeps destructive actions explicit
\item
  lowers to action labels and availability facts
\end{itemize}

\subsection{SliderControl}\label{record-128-slidercontrol}

Labeled numeric control with value, range, and validation semantics.

Variants:

\begin{itemize}
\tightlist
\item
  simple slider
\item
  stepped slider
\item
  range slider
\item
  setting row slider
\end{itemize}

Slots:

\begin{itemize}
\tightlist
\item
  label
\item
  value
\item
  minimum
\item
  maximum
\item
  step
\item
  help text
\item
  validation state
\end{itemize}

Behavior:

\begin{itemize}
\tightlist
\item
  exposes value and range as semantic facts
\item
  supports keyboard and pointer adjustment
\item
  lowers to \texttt{l\allowbreak{}a\allowbreak{}b\allowbreak{}e\allowbreak{}l\allowbreak{}=\allowbreak{}v\allowbreak{}a\allowbreak{}l\allowbreak{}u\allowbreak{}e} with min, max, and step facts
\end{itemize}

\subsection{Grid}\label{record-128-grid}

General cell or item grid for selectable and scannable collections.

Variants:

\begin{itemize}
\tightlist
\item
  fixed grid
\item
  responsive grid
\item
  selectable grid
\item
  virtualized grid
\end{itemize}

Slots:

\begin{itemize}
\tightlist
\item
  items
\item
  selected item
\item
  empty state
\item
  sort/filter controls
\item
  pagination or viewport facts
\end{itemize}

Behavior:

\begin{itemize}
\tightlist
\item
  preserves source order and selection facts
\item
  can host \texttt{A\allowbreak{}p\allowbreak{}p\allowbreak{}l\allowbreak{}i\allowbreak{}c\allowbreak{}a\allowbreak{}t\allowbreak{}i\allowbreak{}o\allowbreak{}n\allowbreak{}C\allowbreak{}a\allowbreak{}r\allowbreak{}d} or other item blocks
\item
  lowers to list/table shape depending on mode
\end{itemize}

\subsection{ApplicationCard}\label{record-128-applicationcard}

Repeated app, tool, document, package, or integration summary.

Variants:

\begin{itemize}
\tightlist
\item
  compact card
\item
  detailed card
\item
  selectable card
\item
  actionable card
\end{itemize}

Slots:

\begin{itemize}
\tightlist
\item
  title
\item
  summary
\item
  status
\item
  metadata
\item
  primary action
\item
  secondary actions
\end{itemize}

Behavior:

\begin{itemize}
\tightlist
\item
  exposes title, status, metadata, and actions semantically
\item
  avoids making visual card shape the only grouping signal
\item
  can live in \texttt{A\allowbreak{}p\allowbreak{}p\allowbreak{}l\allowbreak{}i\allowbreak{}c\allowbreak{}a\allowbreak{}t\allowbreak{}i\allowbreak{}o\allowbreak{}n\allowbreak{}C\allowbreak{}a\allowbreak{}r\allowbreak{}d\allowbreak{}G\allowbreak{}r\allowbreak{}i\allowbreak{}d}, search results, or dashboards
\end{itemize}

\subsection{ApplicationCardGrid}\label{record-128-applicationcardgrid}

Grid of application cards with filtering, selection, and status.

Variants:

\begin{itemize}
\tightlist
\item
  app launcher
\item
  integration gallery
\item
  docs card grid
\item
  selectable dashboard
\end{itemize}

Slots:

\begin{itemize}
\tightlist
\item
  cards
\item
  filter
\item
  sort
\item
  selected card
\item
  empty state
\end{itemize}

Behavior:

\begin{itemize}
\tightlist
\item
  composes \texttt{G\allowbreak{}r\allowbreak{}i\allowbreak{}d} and \texttt{A\allowbreak{}p\allowbreak{}p\allowbreak{}l\allowbreak{}i\allowbreak{}c\allowbreak{}a\allowbreak{}t\allowbreak{}i\allowbreak{}o\allowbreak{}n\allowbreak{}C\allowbreak{}a\allowbreak{}r\allowbreak{}d}
\item
  keeps selected card and active filter visible to agents
\item
  lowers to ordered card summaries
\end{itemize}

\subsection{DataVisualizationPanel}\label{record-128-datavisualizationpanel}

Chart or graph panel with accessible facts and fallback summaries.

Variants:

\begin{itemize}
\tightlist
\item
  sparkline panel
\item
  bar chart panel
\item
  timeline panel
\item
  graph summary panel
\end{itemize}

Slots:

\begin{itemize}
\tightlist
\item
  title
\item
  visualization surface
\item
  legend
\item
  data summary
\item
  selected datum
\item
  actions
\end{itemize}

Behavior:

\begin{itemize}
\tightlist
\item
  exposes chart type, labels, values, and selected datum
\item
  lowers to data summaries when visual rendering is unavailable
\item
  does not make color or geometry the only carrier of meaning
\end{itemize}

\subsection{Carousel}\label{record-128-carousel}

Ordered slide or item viewer with optional animated transitions.

Variants:

\begin{itemize}
\tightlist
\item
  static carousel
\item
  spring slide carousel
\item
  paged carousel
\item
  autoplay carousel
\end{itemize}

Config:

\begin{itemize}
\tightlist
\item
  active index
\item
  transition policy
\item
  spring preset
\item
  autoplay policy
\item
  motion preference
\end{itemize}

Slots:

\begin{itemize}
\tightlist
\item
  slides
\item
  active slide
\item
  previous action
\item
  next action
\item
  page indicators
\item
  pause action
\end{itemize}

Behavior:

\begin{itemize}
\tightlist
\item
  exposes active slide, count, and navigation actions
\item
  respects reduced-motion policy
\item
  lowers to active slide plus ordered slide summaries
\item
  treats spring animation as presentation, not semantic state
\end{itemize}

\subsection{CommentsThread}\label{record-128-commentsthread}

Conversation, review, or annotation surface.

Variants:

\begin{itemize}
\tightlist
\item
  flat comments
\item
  threaded comments
\item
  review comments
\item
  composer with preview
\end{itemize}

Slots:

\begin{itemize}
\tightlist
\item
  comments
\item
  focused comment
\item
  author info
\item
  timestamps
\item
  reactions
\item
  composer
\item
  moderation actions
\end{itemize}

Behavior:

\begin{itemize}
\tightlist
\item
  preserves author, time, reply, and moderation facts
\item
  can compose \texttt{U\allowbreak{}s\allowbreak{}e\allowbreak{}r\allowbreak{}I\allowbreak{}n\allowbreak{}f\allowbreak{}o}, \texttt{F\allowbreak{}o\allowbreak{}r\allowbreak{}m\allowbreak{}P\allowbreak{}a\allowbreak{}n\allowbreak{}e\allowbreak{}l}, and \texttt{B\allowbreak{}u\allowbreak{}t\allowbreak{}t\allowbreak{}o\allowbreak{}n\allowbreak{}S\allowbreak{}e\allowbreak{}t}
\item
  lowers to source-order conversation text
\end{itemize}

\section{Application Domain Candidate Blocks}\label{record-128-application-domain-candidate-blocks}

These blocks are broad app-building pieces. They should usually be cataloged before implementation, then proven by the app family that needs them first.

\subsection{DashboardShell}\label{record-128-dashboardshell}

Operational dashboard frame for metrics, feeds, actions, and drill-down panels.

Common variants:

\begin{itemize}
\tightlist
\item
  executive dashboard
\item
  operations dashboard
\item
  product analytics dashboard
\item
  personal workspace dashboard
\end{itemize}

Composes:

\begin{itemize}
\tightlist
\item
  \texttt{H\allowbreak{}e\allowbreak{}a\allowbreak{}d\allowbreak{}e\allowbreak{}r\allowbreak{}B\allowbreak{}a\allowbreak{}r}
\item
  \texttt{N\allowbreak{}a\allowbreak{}v\allowbreak{}i\allowbreak{}g\allowbreak{}a\allowbreak{}t\allowbreak{}i\allowbreak{}o\allowbreak{}n\allowbreak{}B\allowbreak{}a\allowbreak{}r}
\item
  \texttt{S\allowbreak{}t\allowbreak{}a\allowbreak{}t\allowbreak{}s\allowbreak{}P\allowbreak{}a\allowbreak{}n\allowbreak{}e\allowbreak{}l}
\item
  \texttt{D\allowbreak{}a\allowbreak{}t\allowbreak{}a\allowbreak{}V\allowbreak{}i\allowbreak{}s\allowbreak{}u\allowbreak{}a\allowbreak{}l\allowbreak{}i\allowbreak{}z\allowbreak{}a\allowbreak{}t\allowbreak{}i\allowbreak{}o\allowbreak{}n\allowbreak{}P\allowbreak{}a\allowbreak{}n\allowbreak{}e\allowbreak{}l}
\item
  \texttt{A\allowbreak{}c\allowbreak{}t\allowbreak{}i\allowbreak{}v\allowbreak{}i\allowbreak{}t\allowbreak{}y\allowbreak{}L\allowbreak{}o\allowbreak{}g}
\item
  \texttt{L\allowbreak{}e\allowbreak{}a\allowbreak{}d\allowbreak{}e\allowbreak{}r\allowbreak{}b\allowbreak{}o\allowbreak{}a\allowbreak{}r\allowbreak{}d}
\item
  \texttt{D\allowbreak{}e\allowbreak{}t\allowbreak{}a\allowbreak{}i\allowbreak{}l\allowbreak{}P\allowbreak{}a\allowbreak{}n\allowbreak{}e}
\end{itemize}

\subsection{NavigationBar}\label{record-128-navigationbar}

Horizontal or compact navigation chrome for app sections and top-level routes.

Common variants:

\begin{itemize}
\tightlist
\item
  top navbar
\item
  compact navbar
\item
  tab navbar
\item
  account navbar
\end{itemize}

Composes:

\begin{itemize}
\tightlist
\item
  \texttt{H\allowbreak{}e\allowbreak{}a\allowbreak{}d\allowbreak{}e\allowbreak{}r\allowbreak{}B\allowbreak{}a\allowbreak{}r}
\item
  \texttt{U\allowbreak{}s\allowbreak{}e\allowbreak{}r\allowbreak{}I\allowbreak{}n\allowbreak{}f\allowbreak{}o}
\item
  \texttt{C\allowbreak{}o\allowbreak{}m\allowbreak{}m\allowbreak{}a\allowbreak{}n\allowbreak{}d\allowbreak{}P\allowbreak{}a\allowbreak{}l\allowbreak{}e\allowbreak{}t\allowbreak{}t\allowbreak{}e}
\item
  \texttt{B\allowbreak{}u\allowbreak{}t\allowbreak{}t\allowbreak{}o\allowbreak{}n\allowbreak{}S\allowbreak{}e\allowbreak{}t}
\end{itemize}

\subsection{FeatureBlock}\label{record-128-featureblock}

Single feature item with title, support copy, marker, media, and action.

Common variants:

\begin{itemize}
\tightlist
\item
  compact feature
\item
  illustrated feature
\item
  checklist feature
\item
  feature with proof point
\end{itemize}

Composes:

\begin{itemize}
\tightlist
\item
  \texttt{F\allowbreak{}e\allowbreak{}a\allowbreak{}t\allowbreak{}u\allowbreak{}r\allowbreak{}e\allowbreak{}s\allowbreak{}S\allowbreak{}e\allowbreak{}c\allowbreak{}t\allowbreak{}i\allowbreak{}o\allowbreak{}n}
\item
  \texttt{A\allowbreak{}p\allowbreak{}p\allowbreak{}l\allowbreak{}i\allowbreak{}c\allowbreak{}a\allowbreak{}t\allowbreak{}i\allowbreak{}o\allowbreak{}n\allowbreak{}C\allowbreak{}a\allowbreak{}r\allowbreak{}d}
\item
  \texttt{C\allowbreak{}a\allowbreak{}l\allowbreak{}l\allowbreak{}T\allowbreak{}o\allowbreak{}A\allowbreak{}c\allowbreak{}t\allowbreak{}i\allowbreak{}o\allowbreak{}n}
\end{itemize}

\subsection{CallToAction}\label{record-128-calltoaction}

Focused conversion block for a primary next step.

Common variants:

\begin{itemize}
\tightlist
\item
  inline CTA
\item
  banner CTA
\item
  modal CTA
\item
  footer CTA
\end{itemize}

Composes:

\begin{itemize}
\tightlist
\item
  \texttt{B\allowbreak{}u\allowbreak{}t\allowbreak{}t\allowbreak{}o\allowbreak{}n\allowbreak{}S\allowbreak{}e\allowbreak{}t}
\item
  \texttt{B\allowbreak{}a\allowbreak{}n\allowbreak{}n\allowbreak{}e\allowbreak{}r}
\item
  \texttt{O\allowbreak{}f\allowbreak{}f\allowbreak{}e\allowbreak{}r\allowbreak{}M\allowbreak{}o\allowbreak{}d\allowbreak{}a\allowbreak{}l}
\end{itemize}

\subsection{PricingTiers}\label{record-128-pricingtiers}

Plan comparison and purchase entry block.

Common variants:

\begin{itemize}
\tightlist
\item
  simple tiers
\item
  feature matrix tiers
\item
  usage-based tiers
\item
  enterprise contact tier
\end{itemize}

Composes:

\begin{itemize}
\tightlist
\item
  \texttt{C\allowbreak{}o\allowbreak{}m\allowbreak{}p\allowbreak{}a\allowbreak{}r\allowbreak{}i\allowbreak{}s\allowbreak{}o\allowbreak{}n\allowbreak{}P\allowbreak{}a\allowbreak{}n\allowbreak{}e\allowbreak{}l}
\item
  \texttt{P\allowbreak{}r\allowbreak{}o\allowbreak{}d\allowbreak{}u\allowbreak{}c\allowbreak{}t\allowbreak{}S\allowbreak{}p\allowbreak{}e\allowbreak{}c\allowbreak{}s\allowbreak{}S\allowbreak{}h\allowbreak{}e\allowbreak{}e\allowbreak{}t}
\item
  \texttt{C\allowbreak{}h\allowbreak{}e\allowbreak{}c\allowbreak{}k\allowbreak{}o\allowbreak{}u\allowbreak{}t\allowbreak{}F\allowbreak{}l\allowbreak{}o\allowbreak{}w}
\item
  \texttt{C\allowbreak{}a\allowbreak{}l\allowbreak{}l\allowbreak{}T\allowbreak{}o\allowbreak{}A\allowbreak{}c\allowbreak{}t\allowbreak{}i\allowbreak{}o\allowbreak{}n}
\end{itemize}

\subsection{StatsPanel}\label{record-128-statspanel}

Metric summary block for counts, deltas, health, and status.

Common variants:

\begin{itemize}
\tightlist
\item
  KPI strip
\item
  metric cards
\item
  trend stats
\item
  budget stats
\end{itemize}

Composes:

\begin{itemize}
\tightlist
\item
  \texttt{D\allowbreak{}a\allowbreak{}t\allowbreak{}a\allowbreak{}V\allowbreak{}i\allowbreak{}s\allowbreak{}u\allowbreak{}a\allowbreak{}l\allowbreak{}i\allowbreak{}z\allowbreak{}a\allowbreak{}t\allowbreak{}i\allowbreak{}o\allowbreak{}n\allowbreak{}P\allowbreak{}a\allowbreak{}n\allowbreak{}e\allowbreak{}l}
\item
  \texttt{P\allowbreak{}e\allowbreak{}r\allowbreak{}f\allowbreak{}H\allowbreak{}u\allowbreak{}d}
\item
  \texttt{D\allowbreak{}a\allowbreak{}s\allowbreak{}h\allowbreak{}b\allowbreak{}o\allowbreak{}a\allowbreak{}r\allowbreak{}d\allowbreak{}S\allowbreak{}h\allowbreak{}e\allowbreak{}l\allowbreak{}l}
\end{itemize}

\subsection{ComparisonPanel}\label{record-128-comparisonpanel}

Side-by-side comparison block for products, options, plans, or entities.

Common variants:

\begin{itemize}
\tightlist
\item
  compare A vs B
\item
  compare A vs B vs C
\item
  feature matrix
\item
  ranked comparison
\end{itemize}

Composes:

\begin{itemize}
\tightlist
\item
  \texttt{G\allowbreak{}r\allowbreak{}i\allowbreak{}d}
\item
  \texttt{P\allowbreak{}r\allowbreak{}o\allowbreak{}d\allowbreak{}u\allowbreak{}c\allowbreak{}t\allowbreak{}S\allowbreak{}p\allowbreak{}e\allowbreak{}c\allowbreak{}s\allowbreak{}S\allowbreak{}h\allowbreak{}e\allowbreak{}e\allowbreak{}t}
\item
  \texttt{P\allowbreak{}r\allowbreak{}i\allowbreak{}c\allowbreak{}i\allowbreak{}n\allowbreak{}g\allowbreak{}T\allowbreak{}i\allowbreak{}e\allowbreak{}r\allowbreak{}s}
\end{itemize}

\subsection{Timeline}\label{record-128-timeline}

Chronological event, release, activity, or process surface.

Common variants:

\begin{itemize}
\tightlist
\item
  release timeline
\item
  activity timeline
\item
  process timeline
\item
  audit timeline
\end{itemize}

Composes:

\begin{itemize}
\tightlist
\item
  \texttt{A\allowbreak{}c\allowbreak{}t\allowbreak{}i\allowbreak{}v\allowbreak{}i\allowbreak{}t\allowbreak{}y\allowbreak{}L\allowbreak{}o\allowbreak{}g}
\item
  \texttt{U\allowbreak{}s\allowbreak{}e\allowbreak{}r\allowbreak{}I\allowbreak{}n\allowbreak{}f\allowbreak{}o}
\item
  \texttt{C\allowbreak{}o\allowbreak{}m\allowbreak{}m\allowbreak{}e\allowbreak{}n\allowbreak{}t\allowbreak{}s\allowbreak{}T\allowbreak{}h\allowbreak{}r\allowbreak{}e\allowbreak{}a\allowbreak{}d}
\end{itemize}

\subsection{ProductList}\label{record-128-productlist}

Browsable product collection with filters, sort, status, and actions.

Common variants:

\begin{itemize}
\tightlist
\item
  catalog list
\item
  marketplace grid
\item
  inventory list
\item
  selectable product picker
\end{itemize}

Composes:

\begin{itemize}
\tightlist
\item
  \texttt{G\allowbreak{}r\allowbreak{}i\allowbreak{}d}
\item
  \texttt{P\allowbreak{}r\allowbreak{}o\allowbreak{}d\allowbreak{}u\allowbreak{}c\allowbreak{}t\allowbreak{}C\allowbreak{}a\allowbreak{}t\allowbreak{}e\allowbreak{}g\allowbreak{}o\allowbreak{}r\allowbreak{}y}
\item
  \texttt{A\allowbreak{}p\allowbreak{}p\allowbreak{}l\allowbreak{}i\allowbreak{}c\allowbreak{}a\allowbreak{}t\allowbreak{}i\allowbreak{}o\allowbreak{}n\allowbreak{}C\allowbreak{}a\allowbreak{}r\allowbreak{}d}
\item
  \texttt{S\allowbreak{}e\allowbreak{}a\allowbreak{}r\allowbreak{}c\allowbreak{}h\allowbreak{}R\allowbreak{}e\allowbreak{}s\allowbreak{}u\allowbreak{}l\allowbreak{}t\allowbreak{}s\allowbreak{}P\allowbreak{}a\allowbreak{}n\allowbreak{}e}
\end{itemize}

\subsection{ProductDetails}\label{record-128-productdetails}

Detailed product page or panel for title, media, specs, reviews, and purchase actions.

Common variants:

\begin{itemize}
\tightlist
\item
  product page
\item
  product drawer
\item
  product comparison detail
\item
  purchase detail
\end{itemize}

Composes:

\begin{itemize}
\tightlist
\item
  \texttt{H\allowbreak{}e\allowbreak{}r\allowbreak{}o\allowbreak{}H\allowbreak{}e\allowbreak{}a\allowbreak{}d\allowbreak{}e\allowbreak{}r}
\item
  \texttt{P\allowbreak{}r\allowbreak{}o\allowbreak{}d\allowbreak{}u\allowbreak{}c\allowbreak{}t\allowbreak{}S\allowbreak{}p\allowbreak{}e\allowbreak{}c\allowbreak{}s\allowbreak{}S\allowbreak{}h\allowbreak{}e\allowbreak{}e\allowbreak{}t}
\item
  \texttt{R\allowbreak{}e\allowbreak{}v\allowbreak{}i\allowbreak{}e\allowbreak{}w\allowbreak{}s\allowbreak{}B\allowbreak{}l\allowbreak{}o\allowbreak{}c\allowbreak{}k}
\item
  \texttt{S\allowbreak{}h\allowbreak{}o\allowbreak{}p\allowbreak{}p\allowbreak{}i\allowbreak{}n\allowbreak{}g\allowbreak{}C\allowbreak{}a\allowbreak{}r\allowbreak{}t}
\item
  \texttt{C\allowbreak{}a\allowbreak{}l\allowbreak{}l\allowbreak{}T\allowbreak{}o\allowbreak{}A\allowbreak{}c\allowbreak{}t\allowbreak{}i\allowbreak{}o\allowbreak{}n}
\end{itemize}

\subsection{ProductCategory}\label{record-128-productcategory}

Category navigation and product grouping block.

Common variants:

\begin{itemize}
\tightlist
\item
  category list
\item
  category landing
\item
  faceted category
\item
  nested category tree
\end{itemize}

Composes:

\begin{itemize}
\tightlist
\item
  \texttt{N\allowbreak{}a\allowbreak{}v\allowbreak{}i\allowbreak{}g\allowbreak{}a\allowbreak{}t\allowbreak{}i\allowbreak{}o\allowbreak{}n\allowbreak{}L\allowbreak{}i\allowbreak{}s\allowbreak{}t}
\item
  \texttt{P\allowbreak{}r\allowbreak{}o\allowbreak{}d\allowbreak{}u\allowbreak{}c\allowbreak{}t\allowbreak{}L\allowbreak{}i\allowbreak{}s\allowbreak{}t}
\item
  \texttt{G\allowbreak{}r\allowbreak{}i\allowbreak{}d}
\end{itemize}

\subsection{ProductSpecsSheet}\label{record-128-productspecssheet}

Structured product, plan, or object specifications.

Common variants:

\begin{itemize}
\tightlist
\item
  key-value specs
\item
  grouped specs
\item
  comparison specs
\item
  technical sheet
\end{itemize}

Composes:

\begin{itemize}
\tightlist
\item
  \texttt{I\allowbreak{}n\allowbreak{}s\allowbreak{}p\allowbreak{}e\allowbreak{}c\allowbreak{}t\allowbreak{}o\allowbreak{}r\allowbreak{}P\allowbreak{}a\allowbreak{}n\allowbreak{}e\allowbreak{}l}
\item
  \texttt{C\allowbreak{}o\allowbreak{}m\allowbreak{}p\allowbreak{}a\allowbreak{}r\allowbreak{}i\allowbreak{}s\allowbreak{}o\allowbreak{}n\allowbreak{}P\allowbreak{}a\allowbreak{}n\allowbreak{}e\allowbreak{}l}
\item
  \texttt{P\allowbreak{}r\allowbreak{}o\allowbreak{}d\allowbreak{}u\allowbreak{}c\allowbreak{}t\allowbreak{}D\allowbreak{}e\allowbreak{}t\allowbreak{}a\allowbreak{}i\allowbreak{}l\allowbreak{}s}
\end{itemize}

\subsection{ReviewsBlock}\label{record-128-reviewsblock}

Ratings, reviews, testimonials, or feedback surface.

Common variants:

\begin{itemize}
\tightlist
\item
  rating summary
\item
  review list
\item
  review composer
\item
  testimonial carousel
\end{itemize}

Composes:

\begin{itemize}
\tightlist
\item
  \texttt{C\allowbreak{}o\allowbreak{}m\allowbreak{}m\allowbreak{}e\allowbreak{}n\allowbreak{}t\allowbreak{}s\allowbreak{}T\allowbreak{}h\allowbreak{}r\allowbreak{}e\allowbreak{}a\allowbreak{}d}
\item
  \texttt{U\allowbreak{}s\allowbreak{}e\allowbreak{}r\allowbreak{}I\allowbreak{}n\allowbreak{}f\allowbreak{}o}
\item
  \texttt{C\allowbreak{}a\allowbreak{}r\allowbreak{}o\allowbreak{}u\allowbreak{}s\allowbreak{}e\allowbreak{}l}
\item
  \texttt{F\allowbreak{}o\allowbreak{}r\allowbreak{}m\allowbreak{}P\allowbreak{}a\allowbreak{}n\allowbreak{}e\allowbreak{}l}
\end{itemize}

\subsection{ShoppingCart}\label{record-128-shoppingcart}

Cart contents, quantities, totals, discounts, and checkout entry.

Common variants:

\begin{itemize}
\tightlist
\item
  compact cart
\item
  full cart
\item
  cart drawer
\item
  saved cart
\end{itemize}

Composes:

\begin{itemize}
\tightlist
\item
  \texttt{P\allowbreak{}r\allowbreak{}o\allowbreak{}d\allowbreak{}u\allowbreak{}c\allowbreak{}t\allowbreak{}L\allowbreak{}i\allowbreak{}s\allowbreak{}t}
\item
  \texttt{P\allowbreak{}r\allowbreak{}i\allowbreak{}c\allowbreak{}i\allowbreak{}n\allowbreak{}g\allowbreak{}T\allowbreak{}i\allowbreak{}e\allowbreak{}r\allowbreak{}s}
\item
  \texttt{O\allowbreak{}f\allowbreak{}f\allowbreak{}e\allowbreak{}r\allowbreak{}M\allowbreak{}o\allowbreak{}d\allowbreak{}a\allowbreak{}l}
\item
  \texttt{C\allowbreak{}h\allowbreak{}e\allowbreak{}c\allowbreak{}k\allowbreak{}o\allowbreak{}u\allowbreak{}t\allowbreak{}F\allowbreak{}l\allowbreak{}o\allowbreak{}w}
\end{itemize}

\subsection{CheckoutFlow}\label{record-128-checkoutflow}

Purchase flow from cart review through payment and confirmation.

Common variants:

\begin{itemize}
\tightlist
\item
  single-page checkout
\item
  stepped checkout
\item
  invoice checkout
\item
  confirmation summary
\end{itemize}

Composes:

\begin{itemize}
\tightlist
\item
  \texttt{F\allowbreak{}o\allowbreak{}r\allowbreak{}m\allowbreak{}P\allowbreak{}a\allowbreak{}n\allowbreak{}e\allowbreak{}l}
\item
  \texttt{P\allowbreak{}a\allowbreak{}y\allowbreak{}m\allowbreak{}e\allowbreak{}n\allowbreak{}t\allowbreak{}M\allowbreak{}e\allowbreak{}t\allowbreak{}h\allowbreak{}o\allowbreak{}d\allowbreak{}s}
\item
  \texttt{S\allowbreak{}h\allowbreak{}o\allowbreak{}p\allowbreak{}p\allowbreak{}i\allowbreak{}n\allowbreak{}g\allowbreak{}C\allowbreak{}a\allowbreak{}r\allowbreak{}t}
\item
  \texttt{E\allowbreak{}r\allowbreak{}r\allowbreak{}o\allowbreak{}r\allowbreak{}S\allowbreak{}c\allowbreak{}r\allowbreak{}e\allowbreak{}e\allowbreak{}n}
\end{itemize}

\subsection{PaymentMethods}\label{record-128-paymentmethods}

Payment method selection and validation block.

Common variants:

\begin{itemize}
\tightlist
\item
  saved methods
\item
  card entry
\item
  provider handoff
\item
  invoice terms
\end{itemize}

Composes:

\begin{itemize}
\tightlist
\item
  \texttt{F\allowbreak{}o\allowbreak{}r\allowbreak{}m\allowbreak{}P\allowbreak{}a\allowbreak{}n\allowbreak{}e\allowbreak{}l}
\item
  \texttt{B\allowbreak{}u\allowbreak{}t\allowbreak{}t\allowbreak{}o\allowbreak{}n\allowbreak{}S\allowbreak{}e\allowbreak{}t}
\item
  \texttt{C\allowbreak{}h\allowbreak{}e\allowbreak{}c\allowbreak{}k\allowbreak{}o\allowbreak{}u\allowbreak{}t\allowbreak{}F\allowbreak{}l\allowbreak{}o\allowbreak{}w}
\end{itemize}

\subsection{OfferModal}\label{record-128-offermodal}

Promotional, upgrade, discount, or retention offer.

Common variants:

\begin{itemize}
\tightlist
\item
  upgrade offer
\item
  coupon offer
\item
  trial-ending offer
\item
  retention offer
\end{itemize}

Composes:

\begin{itemize}
\tightlist
\item
  \texttt{M\allowbreak{}o\allowbreak{}d\allowbreak{}a\allowbreak{}l\allowbreak{}D\allowbreak{}i\allowbreak{}a\allowbreak{}l\allowbreak{}o\allowbreak{}g\allowbreak{}S\allowbreak{}t\allowbreak{}a\allowbreak{}c\allowbreak{}k}
\item
  \texttt{C\allowbreak{}a\allowbreak{}l\allowbreak{}l\allowbreak{}T\allowbreak{}o\allowbreak{}A\allowbreak{}c\allowbreak{}t\allowbreak{}i\allowbreak{}o\allowbreak{}n}
\item
  \texttt{P\allowbreak{}r\allowbreak{}i\allowbreak{}c\allowbreak{}i\allowbreak{}n\allowbreak{}g\allowbreak{}T\allowbreak{}i\allowbreak{}e\allowbreak{}r\allowbreak{}s}
\end{itemize}

\subsection{BlogGrid}\label{record-128-bloggrid}

Collection of posts, guides, updates, or articles.

Common variants:

\begin{itemize}
\tightlist
\item
  publication grid
\item
  release notes grid
\item
  author archive
\item
  topic archive
\end{itemize}

Composes:

\begin{itemize}
\tightlist
\item
  \texttt{G\allowbreak{}r\allowbreak{}i\allowbreak{}d}
\item
  \texttt{A\allowbreak{}p\allowbreak{}p\allowbreak{}l\allowbreak{}i\allowbreak{}c\allowbreak{}a\allowbreak{}t\allowbreak{}i\allowbreak{}o\allowbreak{}n\allowbreak{}C\allowbreak{}a\allowbreak{}r\allowbreak{}d}
\item
  \texttt{B\allowbreak{}l\allowbreak{}o\allowbreak{}g\allowbreak{}P\allowbreak{}o\allowbreak{}s\allowbreak{}t}
\end{itemize}

\subsection{BlogPost}\label{record-128-blogpost}

Long-form article, update, or guide page.

Common variants:

\begin{itemize}
\tightlist
\item
  article
\item
  release post
\item
  tutorial
\item
  changelog post
\end{itemize}

Composes:

\begin{itemize}
\tightlist
\item
  \texttt{R\allowbreak{}e\allowbreak{}a\allowbreak{}d\allowbreak{}e\allowbreak{}r\allowbreak{}S\allowbreak{}u\allowbreak{}r\allowbreak{}f\allowbreak{}a\allowbreak{}c\allowbreak{}e}
\item
  \texttt{T\allowbreak{}a\allowbreak{}b\allowbreak{}l\allowbreak{}e\allowbreak{}O\allowbreak{}f\allowbreak{}C\allowbreak{}o\allowbreak{}n\allowbreak{}t\allowbreak{}e\allowbreak{}n\allowbreak{}t\allowbreak{}s}
\item
  \texttt{C\allowbreak{}o\allowbreak{}m\allowbreak{}m\allowbreak{}e\allowbreak{}n\allowbreak{}t\allowbreak{}s\allowbreak{}T\allowbreak{}h\allowbreak{}r\allowbreak{}e\allowbreak{}a\allowbreak{}d}
\item
  \texttt{C\allowbreak{}a\allowbreak{}l\allowbreak{}l\allowbreak{}T\allowbreak{}o\allowbreak{}A\allowbreak{}c\allowbreak{}t\allowbreak{}i\allowbreak{}o\allowbreak{}n}
\end{itemize}

\subsection{Leaderboard}\label{record-128-leaderboard}

Ranked list of users, runs, scores, packages, or entities.

Common variants:

\begin{itemize}
\tightlist
\item
  score leaderboard
\item
  performance leaderboard
\item
  contribution leaderboard
\item
  ranked products
\end{itemize}

Composes:

\begin{itemize}
\tightlist
\item
  \texttt{G\allowbreak{}r\allowbreak{}i\allowbreak{}d}
\item
  \texttt{U\allowbreak{}s\allowbreak{}e\allowbreak{}r\allowbreak{}I\allowbreak{}n\allowbreak{}f\allowbreak{}o}
\item
  \texttt{S\allowbreak{}t\allowbreak{}a\allowbreak{}t\allowbreak{}s\allowbreak{}P\allowbreak{}a\allowbreak{}n\allowbreak{}e\allowbreak{}l}
\end{itemize}

\subsection{ChatPanel}\label{record-128-chatpanel}

Conversation surface for human chat, support, agents, or collaboration.

Common variants:

\begin{itemize}
\tightlist
\item
  support chat
\item
  agent chat
\item
  team chat
\item
  command chat
\end{itemize}

Composes:

\begin{itemize}
\tightlist
\item
  \texttt{C\allowbreak{}o\allowbreak{}m\allowbreak{}m\allowbreak{}e\allowbreak{}n\allowbreak{}t\allowbreak{}s\allowbreak{}T\allowbreak{}h\allowbreak{}r\allowbreak{}e\allowbreak{}a\allowbreak{}d}
\item
  \texttt{U\allowbreak{}s\allowbreak{}e\allowbreak{}r\allowbreak{}I\allowbreak{}n\allowbreak{}f\allowbreak{}o}
\item
  \texttt{F\allowbreak{}o\allowbreak{}r\allowbreak{}m\allowbreak{}P\allowbreak{}a\allowbreak{}n\allowbreak{}e\allowbreak{}l}
\item
  \texttt{N\allowbreak{}o\allowbreak{}t\allowbreak{}i\allowbreak{}f\allowbreak{}i\allowbreak{}c\allowbreak{}a\allowbreak{}t\allowbreak{}i\allowbreak{}o\allowbreak{}n\allowbreak{}S\allowbreak{}t\allowbreak{}a\allowbreak{}c\allowbreak{}k}
\end{itemize}

\subsection{TerminalPanel}\label{record-128-terminalpanel}

Embedded terminal, command runner, REPL, log stream, or transcript viewer.

Common variants:

\begin{itemize}
\tightlist
\item
  terminal session
\item
  command runner
\item
  readonly transcript
\item
  log tail
\end{itemize}

Composes:

\begin{itemize}
\tightlist
\item
  \texttt{A\allowbreak{}p\allowbreak{}p\allowbreak{}S\allowbreak{}h\allowbreak{}e\allowbreak{}l\allowbreak{}l}
\item
  \texttt{C\allowbreak{}o\allowbreak{}m\allowbreak{}m\allowbreak{}a\allowbreak{}n\allowbreak{}d\allowbreak{}P\allowbreak{}a\allowbreak{}l\allowbreak{}e\allowbreak{}t\allowbreak{}t\allowbreak{}e}
\item
  \texttt{A\allowbreak{}c\allowbreak{}t\allowbreak{}i\allowbreak{}v\allowbreak{}i\allowbreak{}t\allowbreak{}y\allowbreak{}L\allowbreak{}o\allowbreak{}g}
\item
  \texttt{P\allowbreak{}e\allowbreak{}r\allowbreak{}f\allowbreak{}H\allowbreak{}u\allowbreak{}d}
\end{itemize}

\section{Remaining Useful Candidates}\label{record-128-remaining-useful-candidates}

With the current list, the catalog can cover most app scaffolds. The remaining blocks I would still keep in reserve are:

\begin{itemize}
\tightlist
\item
  OnboardingFlow
\item
  SearchFilterBar
\item
  UploadDropzone
\item
  FilePreview
\item
  CalendarPanel
\item
  KanbanBoard
\item
  MapPanel
\item
  PermissionsMatrix
\item
  TeamMembersPanel
\item
  AuditLog
\item
  ImportExportPanel
\item
  EmptyState
\item
  ReleaseNotesPanel
\item
  AIConversationPanel
\item
  DiffViewer
\end{itemize}

\section{Second-Wave Blocks}\label{record-128-second-wave-blocks}

These are likely useful, but should follow the first catalog once the block contract is proven:

\begin{itemize}
\tightlist
\item
  DashboardShell
\item
  NavigationBar
\item
  HeaderBar
\item
  FooterBar
\item
  TableOfContents
\item
  HeroHeader
\item
  FeatureBlock
\item
  FeaturesSection
\item
  FAQSection
\item
  Banner
\item
  CallToAction
\item
  PricingTiers
\item
  StatsPanel
\item
  ComparisonPanel
\item
  ErrorScreen
\item
  UserInfo
\item
  FormPanel
\item
  ButtonSet
\item
  SliderControl
\item
  Grid
\item
  ApplicationCard
\item
  ApplicationCardGrid
\item
  DataVisualizationPanel
\item
  Carousel
\item
  CommentsThread
\item
  Timeline
\item
  ProductList
\item
  ProductDetails
\item
  ProductCategory
\item
  ProductSpecsSheet
\item
  ReviewsBlock
\item
  ShoppingCart
\item
  CheckoutFlow
\item
  PaymentMethods
\item
  OfferModal
\item
  BlogGrid
\item
  BlogPost
\item
  Leaderboard
\item
  ChatPanel
\item
  TerminalPanel
\item
  WorkbenchTable
\item
  FormFlow
\item
  ActivityLog
\item
  EmptyErrorState
\item
  DetailPane
\item
  SearchResultsPane
\item
  FileExplorerPane
\end{itemize}

\section{First Implementation Set}\label{record-128-first-implementation-set}

The first implementation should be small:

\begin{enumerate}
\def\labelenumi{\arabic{enumi}.}
\tightlist
\item
  AppShell
\item
  ReaderSurface
\item
  InspectorPanel
\end{enumerate}

These prove shell/content/inspectable composition without dragging in auth, workspace docking, or complex table behavior.

\section{Terminal Screen Sketches}\label{record-128-terminal-screen-sketches}

These sketches are not renderer snapshots. They are shape contracts for the kind of TUI surfaces the first blocks should make easy to assemble.

\texttt{T\allowbreak{}i\allowbreak{}t\allowbreak{}l\allowbreak{}e\allowbreak{}S\allowbreak{}c\allowbreak{}r\allowbreak{}e\allowbreak{}e\allowbreak{}n} with a simple press-enter gate:

\begin{CodeBlock}
+------------------------------------------+
|                                          |
|                 BIJOU                    |
|        Terminal apps with runtime truth   |
|                                          |
|              Press Enter                 |
|                                          |
|  v5.0.0                         Help     |
+------------------------------------------+
\end{CodeBlock}

\texttt{T\allowbreak{}i\allowbreak{}t\allowbreak{}l\allowbreak{}e\allowbreak{}S\allowbreak{}c\allowbreak{}r\allowbreak{}e\allowbreak{}e\allowbreak{}n} composing \texttt{L\allowbreak{}o\allowbreak{}a\allowbreak{}d\allowbreak{}i\allowbreak{}n\allowbreak{}g\allowbreak{}S\allowbreak{}c\allowbreak{}r\allowbreak{}e\allowbreak{}e\allowbreak{}n}:

\begin{CodeBlock}
+ Bijou --------------------------------------------------- startup: loading +
|                                                                          |
|                         Preparing workspace                              |
|                                                                          |
|  [############----------------------------]  32%                          |
|                                                                          |
|  stage 2/5: indexing docs                                                |
|                                                                          |
|  Cancel                                                                  |
+ static output reports stage and progress, not animation frames ----------+
\end{CodeBlock}

\texttt{T\allowbreak{}i\allowbreak{}t\allowbreak{}l\allowbreak{}e\allowbreak{}S\allowbreak{}c\allowbreak{}r\allowbreak{}e\allowbreak{}e\allowbreak{}n} with an auth gateway:

\begin{CodeBlock}
+ Bijou -------------------------------------------------- account required +
| Welcome back                                                             |
|                                                                          |
| Email                                                                    |
| [james@example.com                    ]                                  |
|                                                                          |
| Password                                                                 |
| [********************************      ]                                  |
|                                                                          |
| [Sign in]     [Create account]                                  Help     |
+ lower modes expose available actions without exposing credentials --------+
\end{CodeBlock}

Wide \texttt{A\allowbreak{}p\allowbreak{}p\allowbreak{}S\allowbreak{}h\allowbreak{}e\allowbreak{}l\allowbreak{}l} with \texttt{R\allowbreak{}e\allowbreak{}a\allowbreak{}d\allowbreak{}e\allowbreak{}r\allowbreak{}S\allowbreak{}u\allowbreak{}r\allowbreak{}f\allowbreak{}a\allowbreak{}c\allowbreak{}e} and \texttt{I\allowbreak{}n\allowbreak{}s\allowbreak{}p\allowbreak{}e\allowbreak{}c\allowbreak{}t\allowbreak{}o\allowbreak{}r\allowbreak{}P\allowbreak{}a\allowbreak{}n\allowbreak{}e\allowbreak{}l}:

\begin{CodeBlock}
+ Bijou Docs ---------------------------------------------- mode: interactive +
| Nav                     | Reader                         | Inspector        |
| > Guide                 | # Layout envelopes             | node docs.body   |
|   API                   |                                | role viewport    |
|   Examples              | Every visible node resolves    | inline 24..72    |
|                         | a layout envelope before       | block 1..21      |
|                         | rendering.                     | reason stack     |
|                         |                                | actions          |
|                         |                                | - explain        |
|                         |                                | - copy facts     |
+ status: section 2/8 ------------------------------ notifications: 1 pending +
\end{CodeBlock}

Narrow \texttt{A\allowbreak{}p\allowbreak{}p\allowbreak{}S\allowbreak{}h\allowbreak{}e\allowbreak{}l\allowbreak{}l} preserving primary content:

\begin{CodeBlock}
+ Bijou Docs ------------------------------+
| # Layout envelopes                       |
|                                          |
| Every visible node resolves a layout     |
| envelope before rendering.               |
|                                          |
| [nav] [search] [help] [inspect]          |
+ section 2/8 -----------------------------+
\end{CodeBlock}

\texttt{C\allowbreak{}o\allowbreak{}m\allowbreak{}m\allowbreak{}a\allowbreak{}n\allowbreak{}d\allowbreak{}P\allowbreak{}a\allowbreak{}l\allowbreak{}e\allowbreak{}t\allowbreak{}t\allowbreak{}e} hosted by \texttt{M\allowbreak{}o\allowbreak{}d\allowbreak{}a\allowbreak{}l\allowbreak{}D\allowbreak{}i\allowbreak{}a\allowbreak{}l\allowbreak{}o\allowbreak{}g\allowbreak{}S\allowbreak{}t\allowbreak{}a\allowbreak{}c\allowbreak{}k}:

\begin{CodeBlock}
| App content is inert while the modal stack owns focus                  |
|                                                                        |
|   + Command Palette ------------------------------------------------+  |
|   | > layout                                                       |  |
|   |                                                                |  |
|   | Commands                                                       |  |
|   |   Explain current layout node                                  |  |
|   |   Toggle layout inspector                                      |  |
|   |   Copy layout facts                                            |  |
|   |                                                                |  |
|   | Destinations                                                    |  |
|   |   Layout envelopes                                              |  |
|   |   Viewports and overflow                                        |  |
|   + enter run --------------------------------------------- esc close +  |
\end{CodeBlock}

\texttt{L\allowbreak{}o\allowbreak{}g\allowbreak{}i\allowbreak{}n\allowbreak{}F\allowbreak{}o\allowbreak{}r\allowbreak{}m} credential variants:

\begin{CodeBlock}
+ Sign in: manual password ----------------+
| Email                                    |
| [james@example.com                    ]  |
|                                          |
| Password                                 |
| [********************************      ]  |
|                                          |
| [Sign in]                         Help   |
+ static/pipe output must not expose secret+

+ Sign in: one-time code ------------------+
| Email                                    |
| [james@example.com                    ]  |
|                                          |
| Code                                     |
| [123456                              ]   |
|                                          |
| [Verify]                  Resend    Help |
+ accessible output exposes expiry, not code+

+ Sign in: credential provider ------------+
| Account                                  |
| james@example.com                        |
|                                          |
| Provider                                 |
| @git-stunts/vault adapter: available     |
|                                          |
| [Unlock]                         Cancel  |
+ block sees provider status, not secret   +
\end{CodeBlock}

\section{SVG Block Mockups}\label{record-128-svg-block-mockups}

These SVGs are visual mockups of TUI block shapes. They are not layout diagrams, renderer snapshots, or pixel-perfect specs. The point is to make common block compositions easier to scan in docs and reviews.

App structure block mockups

Wide and narrow app shell mockups showing navigation, reader, and inspector blocks. AppShell wide Bijou Docs interactive

NavigationList \textgreater{} Guide API Examples

ReaderSurface \# Layout envelopes Every visible node resolves layout truth before rendering.

InspectorPanel node docs.body role viewport reason stack - explain - copy facts

status: section 2/8 notifications: 1

AppShell narrow Bijou Docs ReaderSurface \# Layout envelopes Primary content stays visible on small terms. nav search help ins section 2/8

Navigation drawer variant Navigation \textgreater{} Current Search Drawer overlays content focus and source order stay explicit

Inspector as modal body InspectorPanel same facts, different host

App structure blocks show the same semantic regions in wide, narrow, drawer, and modal-hosted forms.

Entry and readiness block mockups

Title, hero, and loading screen variants for app entry flows. TitleScreen: press enter BIJOU runtime truth for terminals Press Enter v5.0.0 Help

TitleScreen + LoadingScreen Preparing workspace stage 2/5: indexing docs Cancel

TitleScreen: auth gateway Welcome back Email james@example.com Password ************ Sign in Create account

Hero: docs landing Bijou Guide Learn by running DOGFOOD. Open Browse

LoadingScreen: retryable failure Workspace unavailable retryable failure Retry Quit

TitleScreen: profile chooser Choose profile \textgreater{} James Guest Enter select Esc quit

Entry blocks separate title, hero, readiness, authentication, and profile selection while allowing those surfaces to compose.

Authentication block mockups

Login and sign-up variants covering password, one-time code, provider, and invite flows. LoginForm: manual password Email james@example.com Password ************ Sign in Help

LoginForm: one-time code Email james@example.com Code 123456 expires in 02:00 Verify Resend

LoginForm: credential provider Account james@example.com Provider @git-stunts/vault ready Unlock Cancel

SignUpForm: password signup Email Password + confirm Create account

SignUpForm: invite code Invite code BIJOU-2026 Terms accepted: yes Continue

SignUpForm: provider handoff Create with provider Use workspace SSO Consent required before handoff

Authentication blocks expose credential modes and provider state without putting secret values into metadata, diagnostics, or lower-mode output.

Command and overlay block mockups

Command palette, help, settings, notifications, modal, and quit dialog block forms. CommandPalette with preview \textgreater{} layout Commands Explain layout node Toggle inspector Copy layout facts Preview node docs.body assigned 72x22 enter run esc close

HelpPanel and SettingsPanel HelpPanel Navigation Enter: open /: search ?: help SettingsPanel Theme {[}dark{]} Density {[}cozy{]} Mode {[}auto{]} Apply

NotificationStack Saved workspace Index still running Sync failed: retry

ModalDialogStack Active modal focus is trapped lower app blocked

QuitConfirmationDialog Quit with unsaved work? 2 docs changed Quit Save Cancel

Command, assistance, notification, and modal blocks show focus ownership and available actions as explicit surface facts.

Runtime diagnostic block mockups

Performance HUD variants showing compact overlay, detailed panel, recorder, and budget warning strip forms. PerfHud: compact overlay fps 60 frame 7ms heap 42MB App content stays behind HUD

PerfHud: detailed panel Runtime metrics render 7ms layout 2ms input 1ms heap 42MB

PerfHud: recorder Capture performance Record Mark samples: 240 last spike: render

PerfHud: budget warning strip layout budget exceeded: 18ms \textgreater{} 12ms actions: explain copy metrics hide

Diagnostic blocks are often existing componentry promoted into explicit slots, variants, config, mode-lowering facts, and story captures.

Section and control block mockups

Header, footer, cards, grids, forms, controls, chart, carousel, FAQ, banner, error, table of contents, and comments block forms. HeaderBar + FooterBar Project Atlas Help content region status ready Save

Grid + ApplicationCard Docs Bench MCP TUI

HeroHeader + Banner Runtime docs Learn layout truth New RE-035 draft Open Compare

FormPanel + controls Preferences Density Mode {[}auto{]} Save Cancel

DataVisualizationPanel Frame time selected: render 18ms

Carousel with spring motion Feature tour Slide 2 of 5 \textless{} Prev Next \textgreater{} Pause

TOC + FAQ + Comments TOC: Overview Layout Blocks FAQ: What is a block? Comments: James replied

ErrorScreen Could not load workspace Retry Open logs Quit

FeaturesSection - retained layout facts - mode lowering - inspector-friendly contracts

Section, control, item, and data blocks standardize recurring app content without making primitive buttons or sliders pretend to be full app blocks.

Application domain block mockups

Product, commerce, publishing, dashboard, leaderboard, chat, and terminal block forms. ProductList + ProductCategory Category: Components Runtime DOGFOOD filter sort compare

ProductDetails + specs Bijou Pro status available specs: 12 facts reviews: 4.8 / 5 Add to cart Compare

PricingTiers + CTA Free \$0 Pro \$20 Ent Start trial

ShoppingCart + CheckoutFlow Cart 2 items \$48 payment: card ending 4242 discount: SPRING Checkout Payment methods

BlogGrid + BlogPost RE-035 Blocks selected post TOC comments CTA

DashboardShell + leaderboard Today runs 128 pass 98\% 1. docs smoke 2. mcp docs 3. tui app

ChatPanel James: ship the floor Agent: facts updated message composer\ldots{}

TerminalPanel \$ npm run docs:inventory Current Truth: ok

ComparisonPanel A vs B vs C speed cost support winner: B for this job

Application-domain blocks cover commerce, publishing, dashboards, chat, and terminal workflows by composing the smaller shell, section, data, and control blocks.

\section{Implementation Outline}\label{record-128-implementation-outline}

\begin{enumerate}
\def\labelenumi{\arabic{enumi}.}
\tightlist
\item
  Done: add the human-readable block catalog and keep it independent from low-level layout implementation.
\item
  Done: define a minimal public \texttt{B\allowbreak{}l\allowbreak{}o\allowbreak{}c\allowbreak{}k\allowbreak{}M\allowbreak{}e\allowbreak{}t\allowbreak{}a\allowbreak{}d\allowbreak{}a\allowbreak{}t\allowbreak{}a}, \texttt{B\allowbreak{}l\allowbreak{}o\allowbreak{}c\allowbreak{}k\allowbreak{}D\allowbreak{}e\allowbreak{}f\allowbreak{}i\allowbreak{}n\allowbreak{}i\allowbreak{}t\allowbreak{}i\allowbreak{}o\allowbreak{}n}, and \texttt{B\allowbreak{}l\allowbreak{}o\allowbreak{}c\allowbreak{}k\allowbreak{}P\allowbreak{}a\allowbreak{}c\allowbreak{}k\allowbreak{}a\allowbreak{}g\allowbreak{}e\allowbreak{}M\allowbreak{}a\allowbreak{}n\allowbreak{}i\allowbreak{}f\allowbreak{}e\allowbreak{}s\allowbreak{}t} shape in \texttt{@\allowbreak{}f\allowbreak{}l\allowbreak{}y\allowbreak{}i\allowbreak{}n\allowbreak{}g\allowbreak{}r\allowbreak{}o\allowbreak{}b\allowbreak{}o\allowbreak{}t\allowbreak{}s\allowbreak{}/\allowbreak{}b\allowbreak{}i\allowbreak{}j\allowbreak{}o\allowbreak{}u}.
\item
  Done: add \texttt{d\allowbreak{}e\allowbreak{}f\allowbreak{}i\allowbreak{}n\allowbreak{}e\allowbreak{}B\allowbreak{}l\allowbreak{}o\allowbreak{}c\allowbreak{}k\allowbreak{}(\allowbreak{})} and \texttt{d\allowbreak{}e\allowbreak{}f\allowbreak{}i\allowbreak{}n\allowbreak{}e\allowbreak{}B\allowbreak{}l\allowbreak{}o\allowbreak{}c\allowbreak{}k\allowbreak{}P\allowbreak{}a\allowbreak{}c\allowbreak{}k\allowbreak{}a\allowbreak{}g\allowbreak{}e\allowbreak{}(\allowbreak{})} authoring helpers, validation/report helpers, and compact summaries.
\item
  Done: formalize provider-bound view data requirements in DX-034 before hardening rendered AppShell APIs.
\item
  Done: attach DX-034 view data contracts and command intents to block definitions for discovery.
\item
  Done: add a structural AppShell composition contract for logical block slots and nested data/command introspection without rendering.
\item
  Done: prove active runtime binding collection, provider snapshot invalidation into new immutable frames, and command intent routing before rendered AppShell work begins.
\item
  Done: add adapter-first \texttt{d\allowbreak{}e\allowbreak{}f\allowbreak{}i\allowbreak{}n\allowbreak{}e\allowbreak{}S\allowbreak{}c\allowbreak{}h\allowbreak{}e\allowbreak{}m\allowbreak{}a\allowbreak{}B\allowbreak{}l\allowbreak{}o\allowbreak{}c\allowbreak{}k\allowbreak{}(\allowbreak{})} support without making Zod, provider lifecycle, AppShell rendering, or DOGFOOD proof part of the core schema-bound block contract.
\item
  Done: add first-party block definitions and story declarations for \texttt{A\allowbreak{}p\allowbreak{}p\allowbreak{}S\allowbreak{}h\allowbreak{}e\allowbreak{}l\allowbreak{}l}, \texttt{R\allowbreak{}e\allowbreak{}a\allowbreak{}d\allowbreak{}e\allowbreak{}r\allowbreak{}S\allowbreak{}u\allowbreak{}r\allowbreak{}f\allowbreak{}a\allowbreak{}c\allowbreak{}e}, and \texttt{I\allowbreak{}n\allowbreak{}s\allowbreak{}p\allowbreak{}e\allowbreak{}c\allowbreak{}t\allowbreak{}o\allowbreak{}r\allowbreak{}P\allowbreak{}a\allowbreak{}n\allowbreak{}e\allowbreak{}l} without making them rendered block products.
\item
  Done: replace the first-party definition placeholders with deterministic rendered proof output and stable lowering facts for \texttt{A\allowbreak{}p\allowbreak{}p\allowbreak{}S\allowbreak{}h\allowbreak{}e\allowbreak{}l\allowbreak{}l}, \texttt{R\allowbreak{}e\allowbreak{}a\allowbreak{}d\allowbreak{}e\allowbreak{}r\allowbreak{}S\allowbreak{}u\allowbreak{}r\allowbreak{}f\allowbreak{}a\allowbreak{}c\allowbreak{}e}, and \texttt{I\allowbreak{}n\allowbreak{}s\allowbreak{}p\allowbreak{}e\allowbreak{}c\allowbreak{}t\allowbreak{}o\allowbreak{}r\allowbreak{}P\allowbreak{}a\allowbreak{}n\allowbreak{}e\allowbreak{}l}.
\item
  Done: add explicit block-tree rendering so nested block declarations render child output without making ordinary \texttt{b\allowbreak{}l\allowbreak{}o\allowbreak{}c\allowbreak{}k\allowbreak{}.\allowbreak{}r\allowbreak{}e\allowbreak{}n\allowbreak{}d\allowbreak{}e\allowbreak{}r\allowbreak{}(\allowbreak{})} recursive.
\item
  Done: expand the standard catalog through the explicit DF-031 status/feedback Block slice.
\item
  Next: add an optional Zod schema adapter package or helper.
\item
  Done: capture interactive, static, pipe, and accessible outputs for the first implementation set.
\item
  Done: prove the first three blocks in DOGFOOD before broadening the catalog.
\item
  Next: add catalog-only variant/config metadata for later blocks without implementing those blocks yet.
\item
  Keep modal stacks, notifications, auth forms, animated carousels, complex controls, and workspace-like behavior outside DX-031; each needs an explicit later block cycle with its own tests and Method evidence.
\end{enumerate}

\section{Tests To Write First}\label{record-128-tests-to-write-first}

\begin{itemize}
\tightlist
\item
  Block catalog tests proving the declared first blocks have metadata.
\item
  Public API tests proving community authors can define a block without relying on DOGFOOD internals.
\item
  Public API tests proving a package can declare exported blocks and Bijou peer compatibility.
\item
  Public API tests proving importing a block package does not mutate global runtime state.
\item
  Public API tests proving a schema-bound block validates data before bind or render output is consumed.
\item
  Public API tests proving an adapter can bind valid data into block slots, config, and semantic facts without a hard Zod dependency.
\item
  Public API tests proving invalid schema data becomes validation facts instead of partial rendered UI.
\item
  Tooling tests proving block metadata can be indexed without rendering a terminal surface.
\item
  Tooling tests proving schema facts can be indexed when the schema adapter exposes them.
\item
  Compiler-path tests proving generated block metadata remains explicit and reviewable instead of hidden behind schema inference.
\item
  Block catalog tests proving each declared variant has a stable id, required slots, and semantic facts.
\item
  Block catalog tests proving config options are visible without storing secret values.
\item
  Block catalog tests proving \texttt{T\allowbreak{}i\allowbreak{}t\allowbreak{}l\allowbreak{}e\allowbreak{}S\allowbreak{}c\allowbreak{}r\allowbreak{}e\allowbreak{}e\allowbreak{}n} can declare press-enter, loading, and auth-gateway variants.
\item
  Block catalog tests proving \texttt{L\allowbreak{}o\allowbreak{}a\allowbreak{}d\allowbreak{}i\allowbreak{}n\allowbreak{}g\allowbreak{}S\allowbreak{}c\allowbreak{}r\allowbreak{}e\allowbreak{}e\allowbreak{}n} exposes readiness, stage, and progress facts without depending on animation.
\item
  Block catalog tests proving proto-blocks such as \texttt{P\allowbreak{}e\allowbreak{}r\allowbreak{}f\allowbreak{}H\allowbreak{}u\allowbreak{}d} can be promoted by declaring slots, variants, config, modes, and semantic facts.
\item
  Block catalog tests proving section, control, item, and data blocks declare their scale and mode-lowering facts.
\item
  Block catalog tests proving application-domain blocks declare their composed lower-level blocks instead of re-owning primitive layout behavior.
\item
  Story tests proving each first block has interactive, static, pipe, and accessible captures.
\item
  DOGFOOD tests proving \texttt{A\allowbreak{}p\allowbreak{}p\allowbreak{}S\allowbreak{}h\allowbreak{}e\allowbreak{}l\allowbreak{}l}, \texttt{R\allowbreak{}e\allowbreak{}a\allowbreak{}d\allowbreak{}e\allowbreak{}r\allowbreak{}S\allowbreak{}u\allowbreak{}r\allowbreak{}f\allowbreak{}a\allowbreak{}c\allowbreak{}e}, and \texttt{I\allowbreak{}n\allowbreak{}s\allowbreak{}p\allowbreak{}e\allowbreak{}c\allowbreak{}t\allowbreak{}o\allowbreak{}r\allowbreak{}P\allowbreak{}a\allowbreak{}n\allowbreak{}e\allowbreak{}l} can compose into a real docs page.
\item
  Mode-lowering tests proving semantic facts survive when visual chrome is removed.
\end{itemize}

\section{Acceptance Criteria}\label{record-128-acceptance-criteria}

\begin{itemize}
\tightlist
\item
  The block catalog defines the first block names, purposes, slots, and behavior contracts.
\item
  Blocks are documented as compositions, not new primitives.
\item
  Blocks have a public authoring API.
\item
  Blocks can be distributed as ordinary NPM packages.
\item
  Block packages declare exported blocks, docs, tags, and Bijou compatibility.
\item
  Importing a block package does not register runtime behavior globally.
\item
  Blocks can bind schema-compliant data objects to slots, config, validation, semantic facts, and render input.
\item
  The core schema-bound block API is adapter-first and does not require Zod as a hard dependency.
\item
  GraphQL, database, or local-storage schema compilers can generate draft block metadata, schema adapters, stories, and semantic facts.
\item
  Tooling can discover block slots, variants, modes, semantic facts, and story captures from metadata before rendering.
\item
  Blocks are grouped by scale so primitive components and app-level blocks do not blur together.
\item
  Blocks can declare variants and configurable options.
\item
  Title and loading blocks are separate but composable.
\item
  \texttt{T\allowbreak{}i\allowbreak{}t\allowbreak{}l\allowbreak{}e\allowbreak{}S\allowbreak{}c\allowbreak{}r\allowbreak{}e\allowbreak{}e\allowbreak{}n} can represent press-enter, loading-gate, and auth-gateway entry flows.
\item
  \texttt{L\allowbreak{}o\allowbreak{}a\allowbreak{}d\allowbreak{}i\allowbreak{}n\allowbreak{}g\allowbreak{}S\allowbreak{}c\allowbreak{}r\allowbreak{}e\allowbreak{}e\allowbreak{}n} can represent indeterminate, determinate, staged, and retryable readiness states.
\item
  Existing block-shaped componentry can be promoted without rewriting it.
\item
  \texttt{P\allowbreak{}e\allowbreak{}r\allowbreak{}f\allowbreak{}H\allowbreak{}u\allowbreak{}d} is captured as a runtime diagnostics block candidate.
\item
  Headers, footers, grids, cards, controls, charts, carousels, FAQs, banners, errors, tables of contents, and comments are captured as candidate blocks.
\item
  Product, commerce, publishing, dashboard, leaderboard, chat, and terminal workflows are captured as application-domain candidate blocks.
\item
  Auth blocks can describe manual-password, one-time-code, and credential-provider modes without exposing credential values.
\item
  The original first implementation milestone landed as \texttt{A\allowbreak{}p\allowbreak{}p\allowbreak{}S\allowbreak{}h\allowbreak{}e\allowbreak{}l\allowbreak{}l}, \texttt{R\allowbreak{}e\allowbreak{}a\allowbreak{}d\allowbreak{}e\allowbreak{}r\allowbreak{}S\allowbreak{}u\allowbreak{}r\allowbreak{}f\allowbreak{}a\allowbreak{}c\allowbreak{}e}, and \texttt{I\allowbreak{}n\allowbreak{}s\allowbreak{}p\allowbreak{}e\allowbreak{}c\allowbreak{}t\allowbreak{}o\allowbreak{}r\allowbreak{}P\allowbreak{}a\allowbreak{}n\allowbreak{}e\allowbreak{}l}.
\item
  Follow-on blocks are captured or promoted through explicit issue-backed slices instead of hidden first-implementation scope.
\item
  DOGFOOD is named as the proving surface for block behavior.
\end{itemize}

\section{Risks / Unknowns}\label{record-128-risks-unknowns}

\begin{itemize}
\tightlist
\item
  Blocks can become too magical if they hide layout or input behavior instead of exposing slots and facts.
\item
  Bijou's block ecosystem can make the project feel like an operating environment instead of a small terminal widget library. That is acceptable only if runtime truth and app truth stay separate.
\item
  Third-party block packages add supply-chain and compatibility risk, so the first sharing model should be explicit imports from NPM packages, not remote runtime loading.
\item
  A central block registry may be useful later, but it should index package metadata rather than become required infrastructure.
\item
  Schema-bound blocks can become too magical if schemas are allowed to own presentation. Schemas validate data; blocks own mapping and rendering.
\item
  Making Zod a hard core dependency would make the API less portable. Zod should be the first adapter, not the only possible schema source.
\item
  Schema-to-block compilers can overfit to names and shapes. Generated output should be reviewable suggestions, not uninspectable app authority.
\item
  Auth blocks need careful secret-handling posture before implementation.
\item
  Credential-provider blocks need a narrow host adapter seam so vault integrations stay out of block internals.
\item
  Startup blocks need stable lower-mode facts because loading animation is not meaningful in static, pipe, or accessible output.
\item
  Animated carousels need reduced-motion behavior and stable lower-mode facts before spring transitions become implementation scope.
\item
  Proto-block promotion can reveal older component assumptions that do not yet fit the block contract cleanly.
\item
  Modal and notification stacks depend on retained layout hit-testing work.
\item
  Workspace-style docking should remain outside this cycle.
\end{itemize}

\section{Playback}\label{record-128-playback}

DX-031A Block Contract lands the public, metadata-first slice of this cycle.

What is now true:

\begin{itemize}
\tightlist
\item
  \texttt{@\allowbreak{}f\allowbreak{}l\allowbreak{}y\allowbreak{}i\allowbreak{}n\allowbreak{}g\allowbreak{}r\allowbreak{}o\allowbreak{}b\allowbreak{}o\allowbreak{}t\allowbreak{}s\allowbreak{}/\allowbreak{}b\allowbreak{}i\allowbreak{}j\allowbreak{}o\allowbreak{}u} exports \texttt{B\allowbreak{}l\allowbreak{}o\allowbreak{}c\allowbreak{}k\allowbreak{}M\allowbreak{}e\allowbreak{}t\allowbreak{}a\allowbreak{}d\allowbreak{}a\allowbreak{}t\allowbreak{}a}, \texttt{B\allowbreak{}l\allowbreak{}o\allowbreak{}c\allowbreak{}k\allowbreak{}D\allowbreak{}e\allowbreak{}f\allowbreak{}i\allowbreak{}n\allowbreak{}i\allowbreak{}t\allowbreak{}i\allowbreak{}o\allowbreak{}n}, \texttt{B\allowbreak{}l\allowbreak{}o\allowbreak{}c\allowbreak{}k\allowbreak{}P\allowbreak{}a\allowbreak{}c\allowbreak{}k\allowbreak{}a\allowbreak{}g\allowbreak{}e\allowbreak{}M\allowbreak{}a\allowbreak{}n\allowbreak{}i\allowbreak{}f\allowbreak{}e\allowbreak{}s\allowbreak{}t}, \texttt{d\allowbreak{}e\allowbreak{}f\allowbreak{}i\allowbreak{}n\allowbreak{}e\allowbreak{}B\allowbreak{}l\allowbreak{}o\allowbreak{}c\allowbreak{}k\allowbreak{}(\allowbreak{})}, \texttt{d\allowbreak{}e\allowbreak{}f\allowbreak{}i\allowbreak{}n\allowbreak{}e\allowbreak{}B\allowbreak{}l\allowbreak{}o\allowbreak{}c\allowbreak{}k\allowbreak{}P\allowbreak{}a\allowbreak{}c\allowbreak{}k\allowbreak{}a\allowbreak{}g\allowbreak{}e\allowbreak{}(\allowbreak{})}, \texttt{v\allowbreak{}a\allowbreak{}l\allowbreak{}i\allowbreak{}d\allowbreak{}a\allowbreak{}t\allowbreak{}e\allowbreak{}B\allowbreak{}l\allowbreak{}o\allowbreak{}c\allowbreak{}k\allowbreak{}M\allowbreak{}e\allowbreak{}t\allowbreak{}a\allowbreak{}d\allowbreak{}a\allowbreak{}t\allowbreak{}a\allowbreak{}(\allowbreak{})}, validation/report helpers, and compact summary helpers.
\item
  Block metadata declares package identity, block name, family, scale, modes, docs, slots, variants, config options, composed components, semantic facts, stories, examples, and tags.
\item
  Block package manifests declare version, Bijou peer compatibility, exported blocks, docs, and tags without hidden global registration.
\item
  Public API tests prove authors can define blocks and package manifests without importing DOGFOOD internals.
\item
  The design-system blocks entrypoint now points to the current public contract instead of only naming future candidates.
\end{itemize}

DX-031B Schema-Bound Blocks adds the boundary validation slice.

What is now true:

\begin{itemize}
\tightlist
\item
  \texttt{@\allowbreak{}f\allowbreak{}l\allowbreak{}y\allowbreak{}i\allowbreak{}n\allowbreak{}g\allowbreak{}r\allowbreak{}o\allowbreak{}b\allowbreak{}o\allowbreak{}t\allowbreak{}s\allowbreak{}/\allowbreak{}b\allowbreak{}i\allowbreak{}j\allowbreak{}o\allowbreak{}u} exports \texttt{B\allowbreak{}l\allowbreak{}o\allowbreak{}c\allowbreak{}k\allowbreak{}S\allowbreak{}c\allowbreak{}h\allowbreak{}e\allowbreak{}m\allowbreak{}a\allowbreak{}A\allowbreak{}d\allowbreak{}a\allowbreak{}p\allowbreak{}t\allowbreak{}e\allowbreak{}r}, \texttt{d\allowbreak{}e\allowbreak{}f\allowbreak{}i\allowbreak{}n\allowbreak{}e\allowbreak{}B\allowbreak{}l\allowbreak{}o\allowbreak{}c\allowbreak{}k\allowbreak{}S\allowbreak{}c\allowbreak{}h\allowbreak{}e\allowbreak{}m\allowbreak{}a\allowbreak{}A\allowbreak{}d\allowbreak{}a\allowbreak{}p\allowbreak{}t\allowbreak{}e\allowbreak{}r\allowbreak{}(\allowbreak{})}, \texttt{d\allowbreak{}e\allowbreak{}f\allowbreak{}i\allowbreak{}n\allowbreak{}e\allowbreak{}S\allowbreak{}c\allowbreak{}h\allowbreak{}e\allowbreak{}m\allowbreak{}a\allowbreak{}B\allowbreak{}l\allowbreak{}o\allowbreak{}c\allowbreak{}k\allowbreak{}(\allowbreak{})}, \texttt{p\allowbreak{}a\allowbreak{}r\allowbreak{}s\allowbreak{}e\allowbreak{}B\allowbreak{}l\allowbreak{}o\allowbreak{}c\allowbreak{}k\allowbreak{}S\allowbreak{}c\allowbreak{}h\allowbreak{}e\allowbreak{}m\allowbreak{}a\allowbreak{}(\allowbreak{})}, and \texttt{b\allowbreak{}i\allowbreak{}n\allowbreak{}d\allowbreak{}S\allowbreak{}c\allowbreak{}h\allowbreak{}e\allowbreak{}m\allowbreak{}a\allowbreak{}B\allowbreak{}l\allowbreak{}o\allowbreak{}c\allowbreak{}k\allowbreak{}I\allowbreak{}n\allowbreak{}p\allowbreak{}u\allowbreak{}t\allowbreak{}(\allowbreak{})}.
\item
  Schema-bound blocks wrap an existing \texttt{B\allowbreak{}l\allowbreak{}o\allowbreak{}c\allowbreak{}k\allowbreak{}D\allowbreak{}e\allowbreak{}f\allowbreak{}i\allowbreak{}n\allowbreak{}i\allowbreak{}t\allowbreak{}i\allowbreak{}o\allowbreak{}n}, preserving block metadata, data contracts, command intents, and ordinary render ownership.
\item
  Schema adapters validate unknown input into immutable typed data or immutable deterministic issues.
\item
  Schema adapters snapshot their parse/describe callbacks at definition time.
\item
  Binding schema data produces immutable block render input and optional semantic facts without rendering, subscribing, dispatching, resolving providers, or owning AppShell behavior.
\item
  Bind outputs must be plain objects with supported render-input keys, so class instances, built-ins, and backchannel handles cannot silently normalize into empty input.
\end{itemize}

DX-031C First-Party Block Definitions adds the first-party standard block definition slice.

What is now true:

\begin{itemize}
\tightlist
\item
  \texttt{@\allowbreak{}f\allowbreak{}l\allowbreak{}y\allowbreak{}i\allowbreak{}n\allowbreak{}g\allowbreak{}r\allowbreak{}o\allowbreak{}b\allowbreak{}o\allowbreak{}t\allowbreak{}s\allowbreak{}/\allowbreak{}b\allowbreak{}i\allowbreak{}j\allowbreak{}o\allowbreak{}u} exports \texttt{a\allowbreak{}p\allowbreak{}p\allowbreak{}S\allowbreak{}h\allowbreak{}e\allowbreak{}l\allowbreak{}l\allowbreak{}B\allowbreak{}l\allowbreak{}o\allowbreak{}c\allowbreak{}k}, \texttt{r\allowbreak{}e\allowbreak{}a\allowbreak{}d\allowbreak{}e\allowbreak{}r\allowbreak{}S\allowbreak{}u\allowbreak{}r\allowbreak{}f\allowbreak{}a\allowbreak{}c\allowbreak{}e\allowbreak{}B\allowbreak{}l\allowbreak{}o\allowbreak{}c\allowbreak{}k}, \texttt{i\allowbreak{}n\allowbreak{}s\allowbreak{}p\allowbreak{}e\allowbreak{}c\allowbreak{}t\allowbreak{}o\allowbreak{}r\allowbreak{}P\allowbreak{}a\allowbreak{}n\allowbreak{}e\allowbreak{}l\allowbreak{}B\allowbreak{}l\allowbreak{}o\allowbreak{}c\allowbreak{}k}, \texttt{s\allowbreak{}t\allowbreak{}a\allowbreak{}n\allowbreak{}d\allowbreak{}a\allowbreak{}r\allowbreak{}d\allowbreak{}B\allowbreak{}l\allowbreak{}o\allowbreak{}c\allowbreak{}k\allowbreak{}s}, \texttt{s\allowbreak{}t\allowbreak{}a\allowbreak{}n\allowbreak{}d\allowbreak{}a\allowbreak{}r\allowbreak{}d\allowbreak{}B\allowbreak{}l\allowbreak{}o\allowbreak{}c\allowbreak{}k\allowbreak{}S\allowbreak{}t\allowbreak{}o\allowbreak{}r\allowbreak{}i\allowbreak{}e\allowbreak{}s}, and \texttt{s\allowbreak{}t\allowbreak{}a\allowbreak{}n\allowbreak{}d\allowbreak{}a\allowbreak{}r\allowbreak{}d\allowbreak{}B\allowbreak{}l\allowbreak{}o\allowbreak{}c\allowbreak{}k\allowbreak{}P\allowbreak{}a\allowbreak{}c\allowbreak{}k\allowbreak{}a\allowbreak{}g\allowbreak{}e\allowbreak{}M\allowbreak{}a\allowbreak{}n\allowbreak{}i\allowbreak{}f\allowbreak{}e\allowbreak{}s\allowbreak{}t}.
\item
  The first-party \texttt{A\allowbreak{}p\allowbreak{}p\allowbreak{}S\allowbreak{}h\allowbreak{}e\allowbreak{}l\allowbreak{}l}, \texttt{R\allowbreak{}e\allowbreak{}a\allowbreak{}d\allowbreak{}e\allowbreak{}r\allowbreak{}S\allowbreak{}u\allowbreak{}r\allowbreak{}f\allowbreak{}a\allowbreak{}c\allowbreak{}e}, and \texttt{I\allowbreak{}n\allowbreak{}s\allowbreak{}p\allowbreak{}e\allowbreak{}c\allowbreak{}t\allowbreak{}o\allowbreak{}r\allowbreak{}P\allowbreak{}a\allowbreak{}n\allowbreak{}e\allowbreak{}l} blocks are real \texttt{d\allowbreak{}e\allowbreak{}f\allowbreak{}i\allowbreak{}n\allowbreak{}e\allowbreak{}B\allowbreak{}l\allowbreak{}o\allowbreak{}c\allowbreak{}k\allowbreak{}(\allowbreak{})} outputs with valid metadata.
\item
  \texttt{R\allowbreak{}e\allowbreak{}a\allowbreak{}d\allowbreak{}e\allowbreak{}r\allowbreak{}S\allowbreak{}u\allowbreak{}r\allowbreak{}f\allowbreak{}a\allowbreak{}c\allowbreak{}e} and \texttt{I\allowbreak{}n\allowbreak{}s\allowbreak{}p\allowbreak{}e\allowbreak{}c\allowbreak{}t\allowbreak{}o\allowbreak{}r\allowbreak{}P\allowbreak{}a\allowbreak{}n\allowbreak{}e\allowbreak{}l} declare DX-034 data contracts and command intents.
\item
  Reader and Inspector schema-bound wrappers validate boundary data before producing render input.
\item
  AppShell uses semantic slots: \texttt{n\allowbreak{}a\allowbreak{}v\allowbreak{}i\allowbreak{}g\allowbreak{}a\allowbreak{}t\allowbreak{}i\allowbreak{}o\allowbreak{}n}, \texttt{c\allowbreak{}o\allowbreak{}n\allowbreak{}t\allowbreak{}e\allowbreak{}n\allowbreak{}t}, \texttt{i\allowbreak{}n\allowbreak{}s\allowbreak{}p\allowbreak{}e\allowbreak{}c\allowbreak{}t\allowbreak{}o\allowbreak{}r}, \texttt{s\allowbreak{}t\allowbreak{}a\allowbreak{}t\allowbreak{}u\allowbreak{}s}, and \texttt{o\allowbreak{}v\allowbreak{}e\allowbreak{}r\allowbreak{}l\allowbreak{}a\allowbreak{}y\allowbreak{}s}.
\item
  Definition and composition introspection do not render blocks, subscribe, dispatch, resolve providers, or walk the active runtime tree.
\end{itemize}

Still out of scope after this slice:

\begin{itemize}
\tightlist
\item
  Zod adapters and schema-to-block compilers.
\item
  DOGFOOD block galleries and multi-mode story captures for rendered blocks.
\item
  Modal stacks, notifications, auth blocks, workspace docking, and complex domain-specific block packages.
\end{itemize}

DX-031D DOGFOOD Blocks Section publishes the standard blocks through the project's proving surface.

What is now true:

\begin{itemize}
\tightlist
\item
  DOGFOOD exposes a top-level Blocks section with ``What are Blocks'', authoring, pre-made catalog, surface-block, preview, and lowering pages.
\item
  Standard block preview pages are selectable from side navigation and render live \texttt{A\allowbreak{}p\allowbreak{}p\allowbreak{}S\allowbreak{}h\allowbreak{}e\allowbreak{}l\allowbreak{}l}, \texttt{R\allowbreak{}e\allowbreak{}a\allowbreak{}d\allowbreak{}e\allowbreak{}r\allowbreak{}S\allowbreak{}u\allowbreak{}r\allowbreak{}f\allowbreak{}a\allowbreak{}c\allowbreak{}e}, and \texttt{I\allowbreak{}n\allowbreak{}s\allowbreak{}p\allowbreak{}e\allowbreak{}c\allowbreak{}t\allowbreak{}o\allowbreak{}r\allowbreak{}P\allowbreak{}a\allowbreak{}n\allowbreak{}e\allowbreak{}l} examples instead of raw metadata dumps.
\item
  The preview keeps lower-mode output visible as a compact lowering summary.
\item
  The local \texttt{C\allowbreak{}o\allowbreak{}u\allowbreak{}n\allowbreak{}t\allowbreak{}e\allowbreak{}r\allowbreak{}D\allowbreak{}e\allowbreak{}m\allowbreak{}o\allowbreak{}B\allowbreak{}l\allowbreak{}o\allowbreak{}c\allowbreak{}k} fixture proves a block-shaped command-intent surface can update through DOGFOOD without becoming part of the shipped standard block package.
\item
  DOGFOOD tests assert the Blocks section, preview routing, curated catalog, live examples, lowering posture, and fixture command keys.
\end{itemize}

DX-031E Rendered Standard Block Proof replaces the first-party placeholders with deterministic render output.

What is now true:

\begin{itemize}
\tightlist
\item
  \texttt{A\allowbreak{}p\allowbreak{}p\allowbreak{}S\allowbreak{}h\allowbreak{}e\allowbreak{}l\allowbreak{}l}, \texttt{R\allowbreak{}e\allowbreak{}a\allowbreak{}d\allowbreak{}e\allowbreak{}r\allowbreak{}S\allowbreak{}u\allowbreak{}r\allowbreak{}f\allowbreak{}a\allowbreak{}c\allowbreak{}e}, and \texttt{I\allowbreak{}n\allowbreak{}s\allowbreak{}p\allowbreak{}e\allowbreak{}c\allowbreak{}t\allowbreak{}o\allowbreak{}r\allowbreak{}P\allowbreak{}a\allowbreak{}n\allowbreak{}e\allowbreak{}l} render concrete string output for interactive, static, pipe, and accessible modes.
\item
  \texttt{A\allowbreak{}p\allowbreak{}p\allowbreak{}S\allowbreak{}h\allowbreak{}e\allowbreak{}l\allowbreak{}l} renders semantic region output for \texttt{n\allowbreak{}a\allowbreak{}v\allowbreak{}i\allowbreak{}g\allowbreak{}a\allowbreak{}t\allowbreak{}i\allowbreak{}o\allowbreak{}n}, \texttt{c\allowbreak{}o\allowbreak{}n\allowbreak{}t\allowbreak{}e\allowbreak{}n\allowbreak{}t}, \texttt{i\allowbreak{}n\allowbreak{}s\allowbreak{}p\allowbreak{}e\allowbreak{}c\allowbreak{}t\allowbreak{}o\allowbreak{}r}, \texttt{s\allowbreak{}t\allowbreak{}a\allowbreak{}t\allowbreak{}u\allowbreak{}s}, and \texttt{o\allowbreak{}v\allowbreak{}e\allowbreak{}r\allowbreak{}l\allowbreak{}a\allowbreak{}y\allowbreak{}s}.
\item
  \texttt{R\allowbreak{}e\allowbreak{}a\allowbreak{}d\allowbreak{}e\allowbreak{}r\allowbreak{}S\allowbreak{}u\allowbreak{}r\allowbreak{}f\allowbreak{}a\allowbreak{}c\allowbreak{}e} renders content, navigation, and outline output.
\item
  \texttt{I\allowbreak{}n\allowbreak{}s\allowbreak{}p\allowbreak{}e\allowbreak{}c\allowbreak{}t\allowbreak{}o\allowbreak{}r\allowbreak{}P\allowbreak{}a\allowbreak{}n\allowbreak{}e\allowbreak{}l} renders selection, details, and actions output.
\item
  Render facts stay stable across modes so lowering checks do not need to parse rendered text.
\item
  Render results do not expose provider handles, subscription handles, refresh methods, active-tree traversal, or command dispatch callbacks.
\end{itemize}

Still out of scope after this slice:

\begin{itemize}
\tightlist
\item
  Provider subscriptions and cache policy.
\item
  Active runtime traversal.
\item
  Command dispatch integration.
\item
  Production AppShell behavior.
\item
  Hidden catalog expansion beyond \texttt{A\allowbreak{}p\allowbreak{}p\allowbreak{}S\allowbreak{}h\allowbreak{}e\allowbreak{}l\allowbreak{}l}, \texttt{R\allowbreak{}e\allowbreak{}a\allowbreak{}d\allowbreak{}e\allowbreak{}r\allowbreak{}S\allowbreak{}u\allowbreak{}r\allowbreak{}f\allowbreak{}a\allowbreak{}c\allowbreak{}e}, and \texttt{I\allowbreak{}n\allowbreak{}s\allowbreak{}p\allowbreak{}e\allowbreak{}c\allowbreak{}t\allowbreak{}o\allowbreak{}r\allowbreak{}P\allowbreak{}a\allowbreak{}n\allowbreak{}e\allowbreak{}l}.
\end{itemize}

DF-031 Status Feedback Expansion makes the next catalog step explicit.

What is now true:

\begin{itemize}
\tightlist
\item
  Issues \#220 through \#225 moved into \texttt{v\allowbreak{}6\allowbreak{}.\allowbreak{}0\allowbreak{}.\allowbreak{}0} as a focused status/feedback Block slice.
\item
  The standard status/feedback block names now include \texttt{I\allowbreak{}n\allowbreak{}l\allowbreak{}i\allowbreak{}n\allowbreak{}e\allowbreak{}S\allowbreak{}t\allowbreak{}a\allowbreak{}t\allowbreak{}u\allowbreak{}s\allowbreak{}B\allowbreak{}l\allowbreak{}o\allowbreak{}c\allowbreak{}k}, \texttt{I\allowbreak{}n\allowbreak{}F\allowbreak{}l\allowbreak{}o\allowbreak{}w\allowbreak{}S\allowbreak{}t\allowbreak{}a\allowbreak{}t\allowbreak{}u\allowbreak{}s\allowbreak{}B\allowbreak{}l\allowbreak{}o\allowbreak{}c\allowbreak{}k}, \texttt{T\allowbreak{}r\allowbreak{}a\allowbreak{}n\allowbreak{}s\allowbreak{}i\allowbreak{}e\allowbreak{}n\allowbreak{}t\allowbreak{}O\allowbreak{}v\allowbreak{}e\allowbreak{}r\allowbreak{}l\allowbreak{}a\allowbreak{}y\allowbreak{}B\allowbreak{}l\allowbreak{}o\allowbreak{}c\allowbreak{}k}, \texttt{A\allowbreak{}c\allowbreak{}t\allowbreak{}i\allowbreak{}v\allowbreak{}i\allowbreak{}t\allowbreak{}y\allowbreak{}S\allowbreak{}t\allowbreak{}r\allowbreak{}e\allowbreak{}a\allowbreak{}m\allowbreak{}B\allowbreak{}l\allowbreak{}o\allowbreak{}c\allowbreak{}k}, \texttt{S\allowbreak{}h\allowbreak{}o\allowbreak{}r\allowbreak{}t\allowbreak{}c\allowbreak{}u\allowbreak{}t\allowbreak{}C\allowbreak{}u\allowbreak{}e\allowbreak{}B\allowbreak{}l\allowbreak{}o\allowbreak{}c\allowbreak{}k}, and \texttt{P\allowbreak{}r\allowbreak{}o\allowbreak{}g\allowbreak{}r\allowbreak{}e\allowbreak{}s\allowbreak{}s\allowbreak{}I\allowbreak{}n\allowbreak{}d\allowbreak{}i\allowbreak{}c\allowbreak{}a\allowbreak{}t\allowbreak{}o\allowbreak{}r\allowbreak{}B\allowbreak{}l\allowbreak{}o\allowbreak{}c\allowbreak{}k}.
\item
  The public package barrel exposes the matching lower-camel definitions: \texttt{i\allowbreak{}n\allowbreak{}l\allowbreak{}i\allowbreak{}n\allowbreak{}e\allowbreak{}S\allowbreak{}t\allowbreak{}a\allowbreak{}t\allowbreak{}u\allowbreak{}s\allowbreak{}B\allowbreak{}l\allowbreak{}o\allowbreak{}c\allowbreak{}k}, \texttt{i\allowbreak{}n\allowbreak{}F\allowbreak{}l\allowbreak{}o\allowbreak{}w\allowbreak{}S\allowbreak{}t\allowbreak{}a\allowbreak{}t\allowbreak{}u\allowbreak{}s\allowbreak{}B\allowbreak{}l\allowbreak{}o\allowbreak{}c\allowbreak{}k}, \texttt{t\allowbreak{}r\allowbreak{}a\allowbreak{}n\allowbreak{}s\allowbreak{}i\allowbreak{}e\allowbreak{}n\allowbreak{}t\allowbreak{}O\allowbreak{}v\allowbreak{}e\allowbreak{}r\allowbreak{}l\allowbreak{}a\allowbreak{}y\allowbreak{}B\allowbreak{}l\allowbreak{}o\allowbreak{}c\allowbreak{}k}, \texttt{a\allowbreak{}c\allowbreak{}t\allowbreak{}i\allowbreak{}v\allowbreak{}i\allowbreak{}t\allowbreak{}y\allowbreak{}S\allowbreak{}t\allowbreak{}r\allowbreak{}e\allowbreak{}a\allowbreak{}m\allowbreak{}B\allowbreak{}l\allowbreak{}o\allowbreak{}c\allowbreak{}k}, \texttt{s\allowbreak{}h\allowbreak{}o\allowbreak{}r\allowbreak{}t\allowbreak{}c\allowbreak{}u\allowbreak{}t\allowbreak{}C\allowbreak{}u\allowbreak{}e\allowbreak{}B\allowbreak{}l\allowbreak{}o\allowbreak{}c\allowbreak{}k}, and \texttt{p\allowbreak{}r\allowbreak{}o\allowbreak{}g\allowbreak{}r\allowbreak{}e\allowbreak{}s\allowbreak{}s\allowbreak{}I\allowbreak{}n\allowbreak{}d\allowbreak{}i\allowbreak{}c\allowbreak{}a\allowbreak{}t\allowbreak{}o\allowbreak{}r\allowbreak{}B\allowbreak{}l\allowbreak{}o\allowbreak{}c\allowbreak{}k}, plus matching \texttt{*\allowbreak{}S\allowbreak{}c\allowbreak{}h\allowbreak{}e\allowbreak{}m\allowbreak{}a\allowbreak{}A\allowbreak{}d\allowbreak{}a\allowbreak{}p\allowbreak{}t\allowbreak{}e\allowbreak{}r} and \texttt{*\allowbreak{}S\allowbreak{}c\allowbreak{}h\allowbreak{}e\allowbreak{}m\allowbreak{}a\allowbreak{}B\allowbreak{}l\allowbreak{}o\allowbreak{}c\allowbreak{}k} exports.
\item
  The new Blocks use the same DX-031/DX-034 public contract: validated block metadata, view data contracts, schema-bound adapters, command intents, story ids, deterministic render output, and stable lowering facts.
\item
  Status and feedback block expansion is now explicit issue-backed v6 work through \#220 through \#225, not hidden DX-031 scope.
\item
  YAML/JSON declaration remains a future adapter over the typed Blocks API, not a second runtime IR in this slice.
\end{itemize}

\section{Retrospective}\label{record-128-retrospective}

DX-031 is landed for the \texttt{v\allowbreak{}6\allowbreak{}.\allowbreak{}0\allowbreak{}.\allowbreak{}0} release boundary. The shipped release floor is the standard block set: public block metadata, package manifests, adapter-first schema-bound block contracts, first-party \texttt{A\allowbreak{}p\allowbreak{}p\allowbreak{}S\allowbreak{}h\allowbreak{}e\allowbreak{}l\allowbreak{}l}, \texttt{R\allowbreak{}e\allowbreak{}a\allowbreak{}d\allowbreak{}e\allowbreak{}r\allowbreak{}S\allowbreak{}u\allowbreak{}r\allowbreak{}f\allowbreak{}a\allowbreak{}c\allowbreak{}e}, \texttt{I\allowbreak{}n\allowbreak{}s\allowbreak{}p\allowbreak{}e\allowbreak{}c\allowbreak{}t\allowbreak{}o\allowbreak{}r\allowbreak{}P\allowbreak{}a\allowbreak{}n\allowbreak{}e\allowbreak{}l}, \texttt{I\allowbreak{}n\allowbreak{}l\allowbreak{}i\allowbreak{}n\allowbreak{}e\allowbreak{}S\allowbreak{}t\allowbreak{}a\allowbreak{}t\allowbreak{}u\allowbreak{}s\allowbreak{}B\allowbreak{}l\allowbreak{}o\allowbreak{}c\allowbreak{}k}, \texttt{I\allowbreak{}n\allowbreak{}F\allowbreak{}l\allowbreak{}o\allowbreak{}w\allowbreak{}S\allowbreak{}t\allowbreak{}a\allowbreak{}t\allowbreak{}u\allowbreak{}s\allowbreak{}B\allowbreak{}l\allowbreak{}o\allowbreak{}c\allowbreak{}k}, \texttt{T\allowbreak{}r\allowbreak{}a\allowbreak{}n\allowbreak{}s\allowbreak{}i\allowbreak{}e\allowbreak{}n\allowbreak{}t\allowbreak{}O\allowbreak{}v\allowbreak{}e\allowbreak{}r\allowbreak{}l\allowbreak{}a\allowbreak{}y\allowbreak{}B\allowbreak{}l\allowbreak{}o\allowbreak{}c\allowbreak{}k}, \texttt{A\allowbreak{}c\allowbreak{}t\allowbreak{}i\allowbreak{}v\allowbreak{}i\allowbreak{}t\allowbreak{}y\allowbreak{}S\allowbreak{}t\allowbreak{}r\allowbreak{}e\allowbreak{}a\allowbreak{}m\allowbreak{}B\allowbreak{}l\allowbreak{}o\allowbreak{}c\allowbreak{}k}, \texttt{S\allowbreak{}h\allowbreak{}o\allowbreak{}r\allowbreak{}t\allowbreak{}c\allowbreak{}u\allowbreak{}t\allowbreak{}C\allowbreak{}u\allowbreak{}e\allowbreak{}B\allowbreak{}l\allowbreak{}o\allowbreak{}c\allowbreak{}k}, and \texttt{P\allowbreak{}r\allowbreak{}o\allowbreak{}g\allowbreak{}r\allowbreak{}e\allowbreak{}s\allowbreak{}s\allowbreak{}I\allowbreak{}n\allowbreak{}d\allowbreak{}i\allowbreak{}c\allowbreak{}a\allowbreak{}t\allowbreak{}o\allowbreak{}r\allowbreak{}B\allowbreak{}l\allowbreak{}o\allowbreak{}c\allowbreak{}k} definitions and schema-bound wrappers, deterministic rendered output for interactive, static, pipe, and accessible modes, block-tree rendering, and a DOGFOOD Blocks section that proves the standard set in product context.

DOGFOOD multi-mode proof is present for the release slice through the standard block render tests, DOGFOOD preview tests, lowering posture checks, and story declarations. Broader ``full block availability'' remains intentionally larger than this release issue: optional Zod adapters, schema-to-block compilers, modal stacks, notifications, auth blocks, workspace docking, and broad catalog expansion should land as explicit future issues with their own tests and Method evidence.

